% Options for packages loaded elsewhere
\PassOptionsToPackage{unicode}{hyperref}
\PassOptionsToPackage{hyphens}{url}
%
\documentclass[
]{article}
\usepackage{amsmath,amssymb}
\usepackage{lmodern}
\usepackage{iftex}
\ifPDFTeX
  \usepackage[T1]{fontenc}
  \usepackage[utf8]{inputenc}
  \usepackage{textcomp} % provide euro and other symbols
\else % if luatex or xetex
  \usepackage{unicode-math}
  \defaultfontfeatures{Scale=MatchLowercase}
  \defaultfontfeatures[\rmfamily]{Ligatures=TeX,Scale=1}
  \setmainfont[]{STIX Two Text}
\fi
% Use upquote if available, for straight quotes in verbatim environments
\IfFileExists{upquote.sty}{\usepackage{upquote}}{}
\IfFileExists{microtype.sty}{% use microtype if available
  \usepackage[]{microtype}
  \UseMicrotypeSet[protrusion]{basicmath} % disable protrusion for tt fonts
}{}
\makeatletter
\@ifundefined{KOMAClassName}{% if non-KOMA class
  \IfFileExists{parskip.sty}{%
    \usepackage{parskip}
  }{% else
    \setlength{\parindent}{0pt}
    \setlength{\parskip}{6pt plus 2pt minus 1pt}}
}{% if KOMA class
  \KOMAoptions{parskip=half}}
\makeatother
\usepackage{xcolor}
\IfFileExists{xurl.sty}{\usepackage{xurl}}{} % add URL line breaks if available
\IfFileExists{bookmark.sty}{\usepackage{bookmark}}{\usepackage{hyperref}}
\hypersetup{
  pdftitle={AGEL J020613−011417 --- Drafting Fact Sheet for the Methods \& Results Sections},
  pdfauthor={Compiled for R. Córdova Rosado from project notes, memory, METHODS\_AND\_SYSTEMATICS.md, and TESTS\_AND\_DIAGNOSTICS.md},
  hidelinks,
  pdfcreator={LaTeX via pandoc}}
\urlstyle{same} % disable monospaced font for URLs
\usepackage[margin=1in]{geometry}
\usepackage{color}
\usepackage{fancyvrb}
\newcommand{\VerbBar}{|}
\newcommand{\VERB}{\Verb[commandchars=\\\{\}]}
\DefineVerbatimEnvironment{Highlighting}{Verbatim}{commandchars=\\\{\}}
% Add ',fontsize=\small' for more characters per line
\newenvironment{Shaded}{}{}
\newcommand{\AlertTok}[1]{\textcolor[rgb]{1.00,0.00,0.00}{\textbf{#1}}}
\newcommand{\AnnotationTok}[1]{\textcolor[rgb]{0.38,0.63,0.69}{\textbf{\textit{#1}}}}
\newcommand{\AttributeTok}[1]{\textcolor[rgb]{0.49,0.56,0.16}{#1}}
\newcommand{\BaseNTok}[1]{\textcolor[rgb]{0.25,0.63,0.44}{#1}}
\newcommand{\BuiltInTok}[1]{#1}
\newcommand{\CharTok}[1]{\textcolor[rgb]{0.25,0.44,0.63}{#1}}
\newcommand{\CommentTok}[1]{\textcolor[rgb]{0.38,0.63,0.69}{\textit{#1}}}
\newcommand{\CommentVarTok}[1]{\textcolor[rgb]{0.38,0.63,0.69}{\textbf{\textit{#1}}}}
\newcommand{\ConstantTok}[1]{\textcolor[rgb]{0.53,0.00,0.00}{#1}}
\newcommand{\ControlFlowTok}[1]{\textcolor[rgb]{0.00,0.44,0.13}{\textbf{#1}}}
\newcommand{\DataTypeTok}[1]{\textcolor[rgb]{0.56,0.13,0.00}{#1}}
\newcommand{\DecValTok}[1]{\textcolor[rgb]{0.25,0.63,0.44}{#1}}
\newcommand{\DocumentationTok}[1]{\textcolor[rgb]{0.73,0.13,0.13}{\textit{#1}}}
\newcommand{\ErrorTok}[1]{\textcolor[rgb]{1.00,0.00,0.00}{\textbf{#1}}}
\newcommand{\ExtensionTok}[1]{#1}
\newcommand{\FloatTok}[1]{\textcolor[rgb]{0.25,0.63,0.44}{#1}}
\newcommand{\FunctionTok}[1]{\textcolor[rgb]{0.02,0.16,0.49}{#1}}
\newcommand{\ImportTok}[1]{#1}
\newcommand{\InformationTok}[1]{\textcolor[rgb]{0.38,0.63,0.69}{\textbf{\textit{#1}}}}
\newcommand{\KeywordTok}[1]{\textcolor[rgb]{0.00,0.44,0.13}{\textbf{#1}}}
\newcommand{\NormalTok}[1]{#1}
\newcommand{\OperatorTok}[1]{\textcolor[rgb]{0.40,0.40,0.40}{#1}}
\newcommand{\OtherTok}[1]{\textcolor[rgb]{0.00,0.44,0.13}{#1}}
\newcommand{\PreprocessorTok}[1]{\textcolor[rgb]{0.74,0.48,0.00}{#1}}
\newcommand{\RegionMarkerTok}[1]{#1}
\newcommand{\SpecialCharTok}[1]{\textcolor[rgb]{0.25,0.44,0.63}{#1}}
\newcommand{\SpecialStringTok}[1]{\textcolor[rgb]{0.73,0.40,0.53}{#1}}
\newcommand{\StringTok}[1]{\textcolor[rgb]{0.25,0.44,0.63}{#1}}
\newcommand{\VariableTok}[1]{\textcolor[rgb]{0.10,0.09,0.49}{#1}}
\newcommand{\VerbatimStringTok}[1]{\textcolor[rgb]{0.25,0.44,0.63}{#1}}
\newcommand{\WarningTok}[1]{\textcolor[rgb]{0.38,0.63,0.69}{\textbf{\textit{#1}}}}
\usepackage{longtable,booktabs,array}
\usepackage{calc} % for calculating minipage widths
% Correct order of tables after \paragraph or \subparagraph
\usepackage{etoolbox}
\makeatletter
\patchcmd\longtable{\par}{\if@noskipsec\mbox{}\fi\par}{}{}
\makeatother
% Allow footnotes in longtable head/foot
\IfFileExists{footnotehyper.sty}{\usepackage{footnotehyper}}{\usepackage{footnote}}
\makesavenoteenv{longtable}
\usepackage[normalem]{ulem}
% Avoid problems with \sout in headers with hyperref
\pdfstringdefDisableCommands{\renewcommand{\sout}{}}
\setlength{\emergencystretch}{3em} % prevent overfull lines
\providecommand{\tightlist}{%
  \setlength{\itemsep}{0pt}\setlength{\parskip}{0pt}}
\setcounter{secnumdepth}{-\maxdimen} % remove section numbering
% Header include for TESTS_AND_DIAGNOSTICS pandoc -> xelatex build.
% Goal: prevent overlapping text in wide tables.

\usepackage{array}
\usepackage{etoolbox}

% Tighter column padding so long cells don't overflow
\setlength{\tabcolsep}{3pt}
\renewcommand{\arraystretch}{1.15}

% Shrink table fonts so wide rows fit the page
\AtBeginEnvironment{longtable}{\footnotesize}
\AtBeginEnvironment{tabular}{\footnotesize}
\AtBeginEnvironment{tabularx}{\footnotesize}

% Allow URLs / inline code to break across lines
\usepackage{xurl}
\renewcommand{\UrlBreaks}{\do\/\do\-\do\_\do\.\do\=}

% pandoc's `inline code` becomes \texttt{...} — long paths overflow
% table columns. The Lua filter .pandoc_break_paths.lua rewrites
% these as \BreakablePath{...} with \discretionary{}{}{} break
% opportunities already inserted after / _ - . — so we only need to
% set tt-family and allow sloppy line breaking.
\newcommand{\BreakablePath}[1]{{\ttfamily\sloppy #1}}

% --- Symbol fallback for glyphs missing from STIX Two Text -----------------
% STIX Two Text lacks several common math/symbol glyphs in the OT:scri
% feature set we render with (→, ≤, ≥, ∈, ≈, ↔, ★). Map each one to its
% counterpart in STIX Two Math, which has the full set.
\usepackage{newunicodechar}
\newfontfamily\stixfallback{STIX Two Math}[Scale=MatchLowercase]
\newcommand{\fb}[1]{{\stixfallback #1}}

\newunicodechar{→}{\fb{→}}
\newunicodechar{←}{\fb{←}}
\newunicodechar{↔}{\fb{↔}}
\newunicodechar{≤}{\fb{≤}}
\newunicodechar{≥}{\fb{≥}}
\newunicodechar{∈}{\fb{∈}}
\newunicodechar{≈}{\fb{≈}}
\newunicodechar{★}{\fb{★}}
\newunicodechar{⋆}{\fb{⋆}}
\newunicodechar{≪}{\fb{≪}}
\newunicodechar{≫}{\fb{≫}}
\newunicodechar{≠}{\fb{≠}}
\newunicodechar{≳}{\fb{≳}}
\newunicodechar{≲}{\fb{≲}}
\newunicodechar{⊕}{\fb{⊕}}
\newunicodechar{σ}{\fb{σ}}
\newunicodechar{Σ}{\fb{Σ}}
\newunicodechar{Δ}{\fb{Δ}}
\newunicodechar{α}{\fb{α}}
\newunicodechar{β}{\fb{β}}
\newunicodechar{γ}{\fb{γ}}
\newunicodechar{δ}{\fb{δ}}
\newunicodechar{ε}{\fb{ε}}
\newunicodechar{η}{\fb{η}}
\newunicodechar{λ}{\fb{λ}}
\newunicodechar{μ}{\fb{μ}}
\newunicodechar{π}{\fb{π}}
\newunicodechar{ρ}{\fb{ρ}}
\newunicodechar{τ}{\fb{τ}}
\newunicodechar{φ}{\fb{φ}}
\newunicodechar{χ}{\fb{χ}}
\newunicodechar{²}{\fb{²}}
\newunicodechar{³}{\fb{³}}
\newunicodechar{ω}{\fb{ω}}
\newunicodechar{Ω}{\fb{Ω}}
\newunicodechar{∝}{\fb{∝}}
\newunicodechar{√}{\fb{√}}
\newunicodechar{∞}{\fb{∞}}
\newunicodechar{±}{\fb{±}}
\newunicodechar{×}{\fb{×}}
\newunicodechar{÷}{\fb{÷}}
\newunicodechar{−}{\fb{−}}
\newunicodechar{☉}{\fb{☉}}
\newunicodechar{✓}{\fb{✓}}
\newunicodechar{⚠}{\fb{⚠}}
\ifLuaTeX
  \usepackage{selnolig}  % disable illegal ligatures
\fi

\title{AGEL J020613−011417 --- Drafting Fact Sheet for the Methods \&
Results Sections}
\usepackage{etoolbox}
\makeatletter
\providecommand{\subtitle}[1]{% add subtitle to \maketitle
  \apptocmd{\@title}{\par {\large #1 \par}}{}{}
}
\makeatother
\subtitle{Facts, decisions, and provenance keyed to the paper outline
--- last revised 2026-06-12 (lens-model summary from Ferrami draft;
post-M12)}
\author{Compiled for R. Córdova Rosado from project notes, memory,
METHODS\_AND\_SYSTEMATICS.md, and TESTS\_AND\_DIAGNOSTICS.md}
\date{2026-06-08}

\begin{document}
\maketitle

\hypertarget{how-to-use-this-document}{%
\section{How to use this document}\label{how-to-use-this-document}}

This is a \textbf{drafting reference}, not prose for the paper. Every
bullet is a fact or decision we have already made, with its source file
in \BreakablePath{[brackets]} so you can verify it. Where a number is
superseded, the current value is given and the stale one flagged. The
\textbf{\BreakablePath{<!-\discretionary{}{}{}-\discretionary{}{}{} PV:auto:… -\discretionary{}{}{}-\discretionary{}{}{}>}
blocks} (the headline numbers below + the §3.1 budget table) are
\textbf{generated from
\BreakablePath{results/\discretionary{}{}{}PAPER\_\discretionary{}{}{}VALUES.\discretionary{}{}{}json}}
--- do not hand-edit inside the markers; refresh with
\BreakablePath{python scripts/\discretionary{}{}{}paper\_\discretionary{}{}{}values.\discretionary{}{}{}py -\discretionary{}{}{}-\discretionary{}{}{}render DRAFTING\_\discretionary{}{}{}FACTS\_\discretionary{}{}{}paper\_\discretionary{}{}{}2026-\discretionary{}{}{}05-\discretionary{}{}{}29.\discretionary{}{}{}md}.
\textbf{Three gaps} are flagged up front --- they are not in this repo
and you must source them elsewhere before drafting those paragraphs:

\begin{itemize}
\tightlist
\item
  \textbf{G1 --- Lens modeling (Ferrami et al.) --- RESOLVED 2026-06-12
  (provisional).} Summarised into \textbf{§2.3} (method) and
  \textbf{§3.3} (results) from the \textbf{Ferrami et al.~June-12 draft}
  (\BreakablePath{\~/\discretionary{}{}{}Documents/\discretionary{}{}{}AGEL/\discretionary{}{}{}AGEL\_\discretionary{}{}{}0206\_\discretionary{}{}{}ApJL\_\discretionary{}{}{}Figures/\discretionary{}{}{}SMBH\_\discretionary{}{}{}in\_\discretionary{}{}{}Einstein\_\discretionary{}{}{}Spiral\_\discretionary{}{}{}DESJ0206\_\discretionary{}{}{}0114\_\discretionary{}{}{}June12/\discretionary{}{}{}main.\discretionary{}{}{}md}):
  \BreakablePath{pronto} code, sEPL/cEPL/bEPL halos, point-mass vs
  SPS/NFW perturber test, 15-model evidence table, best-model
  parameters, M•--σ placement. \textbf{All numbers PROVISIONAL} (draft
  stamps them \BreakablePath{PLACEHOLDER UNTIL FINAL RUNS}) and on
  \textbf{Planck-2015 cosmology} (≠ our H₀=70) --- see §2.3 flags.
  Lens-modeling repos remain
  \BreakablePath{\~/\discretionary{}{}{}Documents/\discretionary{}{}{}AGEL/\discretionary{}{}{}202509\_\discretionary{}{}{}DESJ0206\_\discretionary{}{}{}modeling/\discretionary{}{}{}},
  \BreakablePath{20251112\_\discretionary{}{}{}DESJ0206\_\discretionary{}{}{}Pyauto\_\discretionary{}{}{}PRONTO/\discretionary{}{}{}}
  for the underlying runs.
\item
  \textbf{G2 --- X-ray / ``truly quiescent'' statement --- PARTLY
  RESOLVED 2026-06-12.} The Ferrami draft §Data supplies the quiescence
  argument (Chandra DR2 X-ray flux limit 1.8×10⁻¹⁵ erg s⁻¹ cm⁻² →
  L\_X\textless10⁴³ ≈ 10⁻⁴ L\_Edd; LOFAR DR3 radio-quiet; no AGN lines
  in KCWI) --- captured in \textbf{§3.3.6}. Still verify the
  Chandra/LOFAR source values directly before asserting them in
  \emph{our} paper.
\item
  \textbf{G3 --- ``We do not account for peculiar velocities.''} This
  sentence is not currently written in any methods doc. It is a true
  statement about our analysis (we use the line-fit systemic z directly)
  --- just add it; nothing to reconcile.
\end{itemize}

\textbf{Headline numbers} --- \emph{generated from
\BreakablePath{results/\discretionary{}{}{}PAPER\_\discretionary{}{}{}VALUES.\discretionary{}{}{}json}
by
\BreakablePath{scripts/\discretionary{}{}{}paper\_\discretionary{}{}{}values.\discretionary{}{}{}py -\discretionary{}{}{}-\discretionary{}{}{}render};
do not hand-edit inside the markers.}

\begin{itemize}
\tightlist
\item
  σ\_e(\textless R\_e) = \textbf{267.31 ± 12.79 km/s} (asym −12.98 /
  +12.61); often rounded \textbf{267 ± 13 km/s}. Measured at the
  best-mask R\_e=2.097″ aperture; the R\_e-source systematic (±5.13,
  best-mask CoG family) is folded into the budget (§3.1).
\item
  log(M⋆/M☉) = \textbf{11.50 +0.09/−0.14 (stat, 10\%)} ---
  aperture-corrected total at the matched 2 R\_e = 4.19″ elliptical
  aperture (empirical aperture + single-Sérsic model wings;
  GAMA/Taylor+2011 fluxscale). Five estimators: raw \textbf{11.21}
  (empirical) · raw+apcorr \textbf{11.38} · filled \textbf{11.38} ·
  \textbf{total 11.50 (headline)} · Sérsic-total \textbf{11.43}.
  Sérsic-only systematic ±0.19 dex; masking-approach ±0.09 dex. See
  §2.1.1b / §3.2.
\item
  R\_e = \textbf{2.097″ = 15.26 kpc} (best-mask F140W+F200LP CoG; method
  systematic ±0.10″). \textbf{Supersedes 2.305″} (expert mask).
\item
  z\_deflector = \textbf{0.67564}; z\_source = \textbf{1.302}.
\end{itemize}

\textbf{Cosmology (Planck 2015 --- adopted 2026-06-12 to match the
companion lens-modeling paper, Ferrami et al.):} FlatLambdaCDM,
\textbf{H₀ = 67.7 km/s/Mpc, Ω\_m,0 = 0.302}; 1″ ≈ \textbf{7.2764 kpc},
\textbf{D\_A = 1500.9 Mpc}, D\_L = 4214.1 Mpc at z = 0.67564. \emph{(Was
H₀=70, Ω\_m=0.3 → 7.04 kpc / D\_A=1452 Mpc before this date.)}
{[}reference\_hst\_jwst\_data\_properties.md{]}

\begin{quote}
\textbf{Cosmology switch --- impact (2026-06-12).} Matching Planck 2015
changes only the \textbf{distance-scaled} quantities, and only mildly
(D\_L and D\_A both ×1.0331 at our fixed z\_l = 0.67564): -
\textbf{R\_e: 14.76 → 15.26 kpc} (+3.3\%). R\_e in \textbf{arcsec
(2.097″) is unchanged}. - \textbf{log M⋆: 11.47 → 11.50} (+0.0282 dex).
This is an \textbf{exact D\_L² rescale} of the flux-calibrated SED mass
(the M/L is color-driven, hence cosmology-independent), so it is a
uniform additive shift of \emph{every} estimator --- \textbf{no Bagpipes
re-fit needed}. Well inside the ±0.09 (stat) / ±0.19 (Sérsic sys) budget
→ \textbf{not a dominant systematic}. - \textbf{σ\_e = 267.31 ± 12.79
km/s and every angular quantity (θ\_E, R\_e″, masks, apertures) are
IDENTICAL} --- velocities and angles do not depend on cosmology. -
\textbf{Implementation:} the PV:auto headline + §3.1 budget are
re-rendered under Planck 2015
(\BreakablePath{scripts/\discretionary{}{}{}paper\_\discretionary{}{}{}values.\discretionary{}{}{}py}:
KPC=7.2764;
\BreakablePath{DLOGM\_\discretionary{}{}{}COSMO=+0.\discretionary{}{}{}0282}
applied to all M⋆ estimators). \textbf{TODO:} drop
\BreakablePath{DLOGM\_\discretionary{}{}{}COSMO} once
\BreakablePath{results/\discretionary{}{}{}aperture\_\discretionary{}{}{}matched\_\discretionary{}{}{}photometry.\discretionary{}{}{}npz}
is regenerated under Planck 2015 upstream, else it double-counts. -
\textbf{Redshift reconciliation --- TODO (sub-dominant).} Ferrami quotes
z\_l = 0.675 / z\_s = 1.303 but \textbf{takes the redshifts FROM us};
our line-fit \textbf{z\_l = 0.67564, z\_s = 1.30263} (notebook 04) are
canonical and should propagate into their final runs. The
\textasciitilde6×10⁻⁴ / 1×10⁻³ differences shift D\_A by \textless0.1\%
--- not a dominant systematic, left as a to-do.
\end{quote}

\begin{quote}
\textbf{Staleness warning:} older docs
(\BreakablePath{METHODS\_\discretionary{}{}{}AND\_\discretionary{}{}{}SYSTEMATICS.\discretionary{}{}{}md}
2026-05-18 body; pre-2026-06-11 handoffs) print superseded σ\_e numbers
(254.85 / 268.98 / 269.62, totals ±18/±13.27). The current headline is
\textbf{267.31 ± 12.79} (post ``best mask throughout'', 2026-06-11,
R\_e=2.097″). Always quote the current value (canonical source:
\BreakablePath{results/\discretionary{}{}{}PAPER\_\discretionary{}{}{}VALUES.\discretionary{}{}{}json}).
\end{quote}

\begin{center}\rule{0.5\linewidth}{0.5pt}\end{center}

\hypertarget{data-and-observations}{%
\section{§1 Data and observations}\label{data-and-observations}}

\emph{All values below are read directly from the FITS headers of the
analysis files (verified 2026-06-08), not from memory. Target:
\textbf{AGEL J020613−011417} = DES J0206−0114, ICRS RA = 31.55611°, Dec
= −1.23817° (HST
\BreakablePath{RA\_\discretionary{}{}{}TARG}/\BreakablePath{DEC\_\discretionary{}{}{}TARG});
deflector z = 0.67564, source z = 1.302. The field is the massive
cluster \textbf{ACT-CL J0206.2−0114} (ACT DR5, Hilton et al.~2021) ---
the JWST pointing target.}

\begin{longtable}[]{@{}
  >{\raggedright\arraybackslash}p{(\columnwidth - 12\tabcolsep) * \real{0.14}}
  >{\raggedright\arraybackslash}p{(\columnwidth - 12\tabcolsep) * \real{0.14}}
  >{\raggedright\arraybackslash}p{(\columnwidth - 12\tabcolsep) * \real{0.14}}
  >{\raggedright\arraybackslash}p{(\columnwidth - 12\tabcolsep) * \real{0.14}}
  >{\raggedright\arraybackslash}p{(\columnwidth - 12\tabcolsep) * \real{0.14}}
  >{\raggedright\arraybackslash}p{(\columnwidth - 12\tabcolsep) * \real{0.14}}
  >{\raggedright\arraybackslash}p{(\columnwidth - 12\tabcolsep) * \real{0.14}}@{}}
\toprule
Facility & Inst. / mode & Band(s) & Program (PI) & UT date & Exp. &
Pixel scale \\
\midrule
\endhead
Keck II & KCWI --- KCRM red (RL) \textbf{+} KCB blue (BL), Medium
slicer, 2×2 & RL 5625--8941 Å (σ\_e); BL λc 4500 Å & \textbf{U204 + K409
+ U002} (3-night combine) & 2024-12 → 2025-11 & \textbf{≈180 min RED +
≈198 min BLUE} (36 × 300 s RED) & 0.300″/spaxel \\
HST & WFC3/UVIS & F200LP & \textbf{16773 (Glazebrook)} & 2022-07-14 &
600.0 s (NDRIZIM 4) & 0.050″/pix \\
HST & WFC3/IR & F140W & \textbf{16773 (Glazebrook)} & 2022-07-14 & 597.7
s (NDRIZIM 3) & 0.080″/pix \\
JWST & NIRCam SW, module B & F150W2 (CLEAR) & \textbf{05594 (Mahler)} &
2024-08-27 & 1836.0 s & 0.03075″/pix \\
JWST & NIRCam LW, module B & F322W2 (CLEAR) & \textbf{05594 (Mahler)} &
2024-08-27 & 1836.0 s & 0.06301″/pix \\
\bottomrule
\end{longtable}

\hypertarget{keckkcwi-ifu-spectroscopy-the-ux3c3_e-data}{%
\subsubsection{§1.1 Keck/KCWI IFU spectroscopy (the σ\_e
data)}\label{keckkcwi-ifu-spectroscopy-the-ux3c3_e-data}}

\begin{itemize}
\item
  \textbf{Keck II / KCWI}, KCRM \textbf{red} arm, \textbf{RL (Red-Low)}
  grating (\BreakablePath{RGRATNAM}), \textbf{Medium} slicer
  (\BreakablePath{IFUNAM}), \textbf{2×2} binning
  (\BreakablePath{BINNING}/\BreakablePath{CCDSUM}); nod-and-shuffle
  \textbf{off} (\BreakablePath{RNASNAM='Open'}).
\item
  \textbf{The headline cube is a 3-NIGHT, 3-PROGRAM combine} (see
  §2.2.1), NOT a single dataset --- the cube header
  (\BreakablePath{DATE-\discretionary{}{}{}OBS} 2025-11-17,
  \BreakablePath{PROGPI} Jones, \BreakablePath{XPOSURE} 300 s) describes
  only the \textbf{first input frame}, so cite the combine, not the
  header:

  \begin{longtable}[]{@{}llll@{}}
  \toprule
  Night (UT) & Program (PI) & RED (RL) & BLUE (BL) \\
  \midrule
  \endhead
  2024 Dec 29 & U204 & 4 × 300 s & 1 × 1320 s \emph{(excluded by
  reducer)} \\
  2025 Aug 30 & K409 (PI TBD) & 12 × 300 s & 4 × 990 s \\
  2025 Nov 17 & U002 (Jones) & 20 × 300 s & 5 × 1320 s \\
  \bottomrule
  \end{longtable}

  → \textbf{36 × 300 s ≈ 180 min RED + ≈ 198 min BLUE on-source.} The
  NEW
  \textbf{\BreakablePath{\_\discretionary{}{}{}mtwdo\_\discretionary{}{}{}}}
  reduction (Master-Twilight + Dome hybrid flats, K. R. Gupta),
  \BreakablePath{Nov17\_\discretionary{}{}{}2025\_\discretionary{}{}{}DESJ0206\_\discretionary{}{}{}RL\_\discretionary{}{}{}combined\_\discretionary{}{}{}mtwdo\_\discretionary{}{}{}icubes\_\discretionary{}{}{}wcs.\discretionary{}{}{}fits};
  dithered (PA 0°/45°/90°); CR rejection via \textbf{astroscrappy}
  (\BreakablePath{fsmode=median},
  \BreakablePath{psffwhm=2.\discretionary{}{}{}50}).
  \BreakablePath{OBSERVER} = Glazebrook, Gupta, Jones, Kacprzak, Alcorn,
  Rhoades, Barone, Tran, Chen, Vasa. \emph{(K409 PI + Aug-30 DIMM seeing
  still TBD for acknowledgements.)}
\item
  \textbf{Both KCWI arms recorded simultaneously.} The \textbf{red} arm
  (\textbf{KCRM}, RL grating, central λ 7150 Å \BreakablePath{RCWAVE}
  --- the cube below) is the σ\_e dataset. The \textbf{blue} arm
  (\textbf{KCWI/KCB}, \textbf{BL (Blue-Low)} grating
  \BreakablePath{BGRATNAM}, central λ \textbf{4500 Å}
  \BreakablePath{BCWAVE}, \textbf{KBlue} filter
  \BreakablePath{BFILTNAM}, N\&S off) is the ≈198 min BLUE above; it is
  \textbf{not used for σ\_e} (the deflector's Ca H+K\ldots Mg b
  absorption falls in the red) but supplies the \textbf{source-z = 1.302
  cross-check} (Fe II UV multiplet; §2.2.2). \textbf{The blue combined
  cube is not in this repo} (only the RL cube) --- quote the blue arm
  from the observing log, not a local file.
\item
  \textbf{Conditions (Nov-17 frame, header-literal):} airmass
  \textbf{1.165} (\BreakablePath{AIRMASS}); \textbf{guide-star FWHM
  1.271″} (\BreakablePath{GUIDFWHM}) --- the adopted seeing; parallactic
  angle 15.6°, rotator 0°. (Per-night airmass/seeing for the other two
  nights not carried in the combined cube.)
\item
  \textbf{Cube:} 100 × 100 spaxels × 3317 λ-pixels;
  \textbf{0.300″/spaxel} (square; from WCS), \textbf{30″ × 30″} FOV;
  wavelength \textbf{5625--8941 Å} at \textbf{1.0 Å/pix}
  (\BreakablePath{CRVAL3}/\BreakablePath{CD3\_\discretionary{}{}{}3}),
  red central λ 7150 Å (\BreakablePath{RCWAVE}). Spectral resolution
  \textbf{R ≈ 10000} at the fit band → \textbf{σ\_inst ≈ 12.6 km/s}
  (well below σ\_e ≈ 270; the LSF is carried into ppxf).
  {[}reference\_kcwi\_data\_properties.md{]}
\item
  Pointing \BreakablePath{RA}/\BreakablePath{DEC} = 02:06:13.38 /
  −01:14:20.8; target \BreakablePath{TARGRA}/\BreakablePath{TARGDEC} =
  02:06:13.58 / −01:14:19.8.
\end{itemize}

\hypertarget{hstwfc3-imaging-f200lp-f140w}{%
\subsubsection{§1.2 HST/WFC3 imaging (F200LP,
F140W)}\label{hstwfc3-imaging-f200lp-f140w}}

\begin{itemize}
\tightlist
\item
  \textbf{Program 16773, PI Karl Glazebrook}
  (\BreakablePath{PROPOSID}/\BreakablePath{PR\_\discretionary{}{}{}INV\_\discretionary{}{}{}*});
  \BreakablePath{TARGNAME} = DESJ0206-0114; both bands \textbf{UT
  2022-07-14} (\BreakablePath{EXPSTART} MJD 59774.30--59774.31),
  ORIENTAT 111.9°, drizzled, units \textbf{ELECTRONS/S}.
\item
  \textbf{F200LP} --- WFC3/\textbf{UVIS}, aperture UVIS2: EXPTIME
  \textbf{600.0 s}, NDRIZIM 4, NCOMBINE 2, drizzle scale
  \textbf{0.050″/pix}; \BreakablePath{PHOTFLAM} = 5.1851×10⁻²⁰,
  \BreakablePath{PHOTPLAM} = 4923.48 Å → \textbf{ZP\_AB = 27.344}.
\item
  \textbf{F140W} --- WFC3/\textbf{IR}: EXPTIME \textbf{597.7 s}, NDRIZIM
  3, NCOMBINE 3, drizzle scale \textbf{0.080″/pix};
  \BreakablePath{PHOTFLAM} = 1.4829×10⁻²⁰, \BreakablePath{PHOTPLAM} =
  13922.91 Å → \textbf{ZP\_AB = 26.446}.
\item
  AB zeropoints are recomputed per-image from
  \BreakablePath{PHOTFLAM}/\BreakablePath{PHOTPLAM} (never hardcoded);
  \BreakablePath{ZP\_\discretionary{}{}{}AB = −2.\discretionary{}{}{}5·log10(PHOTFLAM) − 21.\discretionary{}{}{}10 − 5·log10(PHOTPLAM) + 18.\discretionary{}{}{}6921}.
  \textbf{Drizzle pixel correlation:} empirical noise ≈1.3--2× the
  CCD-equation value → motivates the 10\% flux floor (§2.1.4).
\end{itemize}

\hypertarget{jwstnircam-imaging-f150w2-f322w2}{%
\subsubsection{§1.3 JWST/NIRCam imaging (F150W2,
F322W2)}\label{jwstnircam-imaging-f150w2-f322w2}}

\begin{itemize}
\tightlist
\item
  \textbf{Program 05594} --- \emph{``JWST Cluster SLICE --- Strong
  Lensing and Cluster Evolution''}, \textbf{PI Guillaume Mahler} (SURVEY
  category); \BreakablePath{TARGPROP} = ACT-CL\_J0206.2−0114 (the
  cluster field; the deflector sits in it). Both bands \textbf{UT
  2024-08-27} (\BreakablePath{DATE-\discretionary{}{}{}BEG}
  06:02--06:40), NIRCam \textbf{module B}, units \textbf{MJy/sr} (i2d),
  NINTS 1 / NGROUPS 4.
\item
  \textbf{F150W2} (SHORT channel, CLEAR pupil):
  \BreakablePath{EFFEXPTM}/\BreakablePath{XPOSURE} \textbf{1836.0 s};
  \BreakablePath{PIXAR\_\discretionary{}{}{}SR} = 2.2222×10⁻¹⁴ sr →
  \textbf{0.03075″/pix} (PIXAR\_A2 = 9.454×10⁻⁴ ″²) → \textbf{ZP\_AB =
  28.033}.
\item
  \textbf{F322W2} (LONG channel, CLEAR pupil):
  \BreakablePath{EFFEXPTM}/\BreakablePath{XPOSURE} \textbf{1836.0 s};
  \BreakablePath{PIXAR\_\discretionary{}{}{}SR} = 9.3307×10⁻¹⁴ sr →
  \textbf{0.06301″/pix} (PIXAR\_A2 = 3.970×10⁻³ ″²) → \textbf{ZP\_AB =
  26.475}.
\item
  JWST AB from area:
  \BreakablePath{ZP\_\discretionary{}{}{}AB = −6.\discretionary{}{}{}10 − 2.\discretionary{}{}{}5·log10(PIXAR\_\discretionary{}{}{}SR)}.
  No PAM correction (drizzled/resampled).
\end{itemize}

\begin{center}\rule{0.5\linewidth}{0.5pt}\end{center}

\hypertarget{methods}{%
\section{§2 Methods}\label{methods}}

\hypertarget{photometric-analysis}{%
\subsection{§2.1 Photometric Analysis}\label{photometric-analysis}}

\hypertarget{masking}{%
\subsubsection{§2.1.1 Masking}\label{masking}}

\begin{itemize}
\tightlist
\item
  \textbf{Use the F200LP
  \BreakablePath{\_\discretionary{}{}{}mask.\discretionary{}{}{}fits} as
  the arc mask --- NOT the F140W one.}
  {[}feedback\_masking\_strategy.md;
  reference\_hst\_jwst\_data\_properties.md; METHODS §II.4{]}

  \begin{itemize}
  \tightlist
  \item
    F200LP mask (2512 HST pixels) is hand-tuned to cover \textbf{the arc
    only} and leaves the \textbf{deflector core unmasked} --- it is the
    original arc-tuned segmentation mask.
  \item
    F140W mask (1121 HST pixels) covers \textbf{the arc PLUS the
    deflector core} (the photometry tool's aperture geometry pulled the
    core in); the two masks have \textbf{zero spatial overlap}.
  \item
    Consequence: F140W is valued for its \emph{image} (sharper Sérsic
    fit), never for its mask. Using the F140W mask as an arc mask would
    destroy the deflector-core signal. \#\#\# §2.1.1b Principled
    (reproducible) arc masking --- NEW 2026-05-29
  \end{itemize}
\end{itemize}

The hand-painted masks are now backed by an objective, reproducible
recipe (so a referee can reproduce the arc selection), validated +
applied to all four bands.
{[}\BreakablePath{NOTES\_\discretionary{}{}{}arc\_\discretionary{}{}{}mask\_\discretionary{}{}{}verification\_\discretionary{}{}{}2026-\discretionary{}{}{}05-\discretionary{}{}{}29.\discretionary{}{}{}md},
\BreakablePath{NOTES\_\discretionary{}{}{}photometry\_\discretionary{}{}{}mask\_\discretionary{}{}{}systematics\_\discretionary{}{}{}2026-\discretionary{}{}{}05-\discretionary{}{}{}29.\discretionary{}{}{}md};
\BreakablePath{scripts/\discretionary{}{}{}arc\_\discretionary{}{}{}mask\_\discretionary{}{}{}verification.\discretionary{}{}{}py},
\BreakablePath{mask\_\discretionary{}{}{}method\_\discretionary{}{}{}comparison.\discretionary{}{}{}py},
\BreakablePath{photometry\_\discretionary{}{}{}systematics.\discretionary{}{}{}py}{]}

\begin{itemize}
\item
  \textbf{Locator = F200LP} (highest contrast on the blue z=1.302
  source). Two independent objective selectors reproduce the F200LP hand
  mask: (i) a color cut m\_F200LP−m\_F140W (arc bluer than the red
  deflector); (ii) a \textbf{Sérsic-residual} mask (subtract a 2D
  deflector model, flag pixels whose positive residual exceeds
  \textbf{k·σ\_sky}, where \textbf{k is the detection threshold in units
  of the sky noise σ\_sky} and σ\_sky is the robust background scatter;
  we adopt \textbf{k = 3}). On F200LP the Sérsic-residual mask
  reproduces the expert photometry to \textbf{0.016 mag}, R\_e to 3\%.
\item
  \textbf{Transfer to other bands by reprojecting the F200LP footprint}
  (WCS + \BreakablePath{map\_\discretionary{}{}{}coordinates}).
  Reproduces the expert JWST aperture mags to \textbf{0.01--0.02 mag}.
\item
  \textbf{WCS registration (NEW 2026-06-10).} A raw WCS reprojection
  leaves a \textbf{0.09--0.18″} residual offset between the F200LP arc
  and the JWST/F140W frames (HST UVIS↔IR and HST↔JWST astrometry; the
  JWST i2d GWCS differs from the FITS-header WCS astropy reads). On the
  thin arc that mis-covers one edge and lets a companion slip past the
  mask. We \textbf{cross-correlation-register} the reprojected F200LP
  arc (and the F140W color layer) to each band's own image near the
  deflector
  (\BreakablePath{photometry\_\discretionary{}{}{}systematics.\discretionary{}{}{}register\_\discretionary{}{}{}shift}),
  which drives the residual offset to \textbf{0.00″} (verified F140W
  0.179→0, F150W2 0.154→0, F322W2 0.089→0). Field-galaxy contamination
  check after registration: \textbf{no unmasked source (\textgreater8σ)
  falls inside the photometry aperture} in any band. \emph{Do not}
  anchor on the F200LP deflector centroid instead --- it is arc-biased
  (the faint blue deflector centroid is pulled \textasciitilde0.26″ by
  the arc), so the arc cross-correlation is the right fiducial.
\item
  \textbf{Deflector model = single-component Sérsic} (per user
  2026-06-10). The deflector is an elliptical, so a single Sérsic is the
  physically-motivated light profile --- we do \textbf{not} fit a
  bulge+disk decomposition (an earlier 2-component model was only a
  numerical kludge to make the residual-based mask growth behave; see
  the mask-evolution note below). The single Sérsic is fit per band with
  masked pixels excluded, capturing \textbf{66--87\%} of the box flux
  (F140W is the weakest, with a localized central residual from the
  absent bulge/disk split, but that residual sits inside the arc-free
  core and does not enter the mask).

  \begin{itemize}
  \tightlist
  \item
    \emph{Implementation gotcha (documented so it isn't reintroduced):}
    \BreakablePath{astropy.\discretionary{}{}{}modeling.\discretionary{}{}{}Sersic2D}'s
    \BreakablePath{amplitude} parameter is \textbf{I(r\_eff), not the
    central peak}. Since I(0)/I(r\_eff)=exp(b\_n)≈2150 at n=4,
    initialising \BreakablePath{amplitude} at the data peak
    over-predicts the centre \textasciitilde2000× and LevMar collapses
    the amplitude to ≈0 (the whole galaxy then reads as positive
    residual → catastrophic over-masking). Initialise
    \BreakablePath{amplitude ≈ peak·exp(−b\_\discretionary{}{}{}n)}
    (\textasciitilde peak·0.05 for n≈3) and \textbf{bound it to
    (peak·1e-3, peak·1.5)}, exactly as
    \BreakablePath{sersic\_\discretionary{}{}{}total\_\discretionary{}{}{}photometry.\discretionary{}{}{}fit\_\discretionary{}{}{}sersic2d}.
  \end{itemize}
\item
  \textbf{IR extends the source footprint --- grown band-specifically
  (HST color gate / deep-JWST morphology guard):} the deeper/sharper
  JWST bands see the arc's fuller extent, so we region-grow the F200LP
  arc seed into the IR. The ``is this pixel source, not deflector
  body?'' test is gated two ways depending on the band's depth (a
  candidate must also be a \textgreater k·σ\_sky residual, within
  R\textless5″, and spatially contiguous with the arc seed --- region
  growing via
  \BreakablePath{scipy.\discretionary{}{}{}ndimage.\discretionary{}{}{}label}):

  \begin{enumerate}
  \def\labelenumi{\arabic{enumi}.}
  \tightlist
  \item
    \textbf{HST bands (F200LP, F140W) --- COLOR GATE:} the pixel must be
    \textbf{bluer than the deflector core} by ≥ \textbf{ΔC = 0.5 mag} in
    \textbf{F200LP−F140W} (\BreakablePath{DCOLOR}), at \textbf{k = 3}
    (\BreakablePath{K\_\discretionary{}{}{}EXT}).
  \item
    \textbf{Deep JWST bands (F150W2, F322W2) --- MORPHOLOGY GUARD:} the
    pixel must be \textbf{source-dominated},
    \BreakablePath{data − sky − model > MORPH\_\discretionary{}{}{}FRAC·model}
    (excess beats the single-Sérsic deflector model → it is OUTSIDE the
    Sérsic-dominated region), at the lower \textbf{k = 2}
    (\BreakablePath{K\_\discretionary{}{}{}EXT\_\discretionary{}{}{}JWST}).
    \textbf{No HST color is required}, so the deep JWST data captures
    the \textbf{diffuse} arc emission HST is too shallow to
    color-confirm. Because the single Sérsic already accounts for
    66--87\% of the deflector flux, \BreakablePath{resid > model}
    rejects the deflector body (model-dominated) while admitting the
    diffuse source (excess-dominated) --- verified: raw mag stable (no
    deflector loss), 0 unmasked \textgreater8σ sources in the aperture,
    the deflector core \textless0.4″ never masked. F150W2 IR-extension
    grows 1049→5288 px (the diffuse arc), F322W2 415→1230 px.
  \end{enumerate}

  \begin{itemize}
  \tightlist
  \item
    \textbf{Why the color gate is needed (and is the right physics).} A
    \emph{single} Sérsic, being the correct elliptical model, still
    under-fits the bright extended IR galaxy body slightly, leaving a
    low-level \textbf{positive} residual over the \emph{red} galaxy.
    Residual-only growth (criterion 1 alone, the old 2-component recipe)
    therefore sweeps red galaxy-body light --- and even neighbouring
    companions --- into the ``arc'' mask (verified: F150W2 mask balloons
    and raw mags fall 1--2 mag). The z=1.302 source is a \textbf{blue,
    star-forming} lensed galaxy; the deflector is a \textbf{red,
    passive} elliptical. Requiring the grown pixels to be \textbf{bluer
    than the deflector core} (criterion 2) exploits this intrinsic
    4000-Å-break contrast to keep the blue source extension while
    rejecting the red body. The mask thus becomes a property of the
    \textbf{source color}, not of the (imperfect) deflector residual.
    This is the same physical contrast used by the validated F200LP
    color selector (i) above; we are simply applying it per-pixel during
    the IR growth.
  \item
    \textbf{Color construction.} Per target-band grid, reproject the
    F200LP and F140W surface brightness
    (\BreakablePath{scripts.\discretionary{}{}{}arc\_\discretionary{}{}{}mask\_\discretionary{}{}{}verification.\discretionary{}{}{}reproject\_\discretionary{}{}{}intensity},
    WCS + bilinear \BreakablePath{map\_\discretionary{}{}{}coordinates})
    and form
    \BreakablePath{−2.\discretionary{}{}{}5·log10(SB\_\discretionary{}{}{}F200LP/\discretionary{}{}{}SB\_\discretionary{}{}{}F140W)}
    (arc → more negative). Constant ZP offsets \textbf{cancel} because
    we threshold on the contrast vs the deflector-core color (median in
    a 0.4--0.8″ annulus), not on an absolute color.
  \item
    \textbf{Caveats.} (a) \emph{(RESOLVED 2026-06-10)} the HST color
    gate is conservative where HST is shallow and was missing faint
    diffuse JWST arc emission (user catch via the deep-stretch
    \BreakablePath{\_\discretionary{}{}{}diffuse\_\discretionary{}{}{}mask\_\discretionary{}{}{}check});
    the \textbf{deep-JWST morphology guard} above now captures it using
    JWST's own sensitivity + the source-dominance test, without
    over-masking the deflector (raw mag stable). (b) ΔC = 0.5 mag and
    MORPH\_FRAC = 1.0 are tunables; the resulting masks reproduce the
    expert hand-mask M⋆ (11.33), our calibration that they neither over-
    nor under-mask. (c) The morphology guard assumes the single Sérsic
    captures the deflector majority --- true here (66--87\%); if a
    future target has a worse single-Sérsic fit, re-check that
    \BreakablePath{resid>model} still protects the body.
  \end{itemize}
\item
  \textbf{raw vs fill-in:} raw photutils (masked) discards the deflector
  light under the arc; the single-Sérsic fill-in (replace masked pixels
  with the model) recovers it. Under the color-gated masks the fill-in
  correction is \textbf{−0.19 to −0.27 mag} per band (smaller than the
  old 2-component +0.18 to +0.96 mag, because the cleaner masks remove
  less deflector light). The \textbf{headline is the empirical raw value
  with the best (per-band, color-gated) mask}; the fill-in is carried as
  a \textbf{one-sided +sys} (deflector light the mask removes).
  \textbf{PSF effect on the fill} stays negligible
  (\BreakablePath{scripts/\discretionary{}{}{}psf\_\discretionary{}{}{}fill\_\discretionary{}{}{}model.\discretionary{}{}{}py}:
  ≤0.004 mag → ΔM⋆ ≪0.01 dex; arc is outside the PSF core).
\item
  \textbf{per-band vs global mask} (union of all per-band masks):
  masking systematic \textbf{±0.086 dex} (10\%); the raw↔filled
  (under-arc) term dominates (±0.057), the per-band↔global
  mask-definition term ±0.028. See §3.2.
\item
  \textbf{Mask-evolution note (why we settled on the single-Sérsic +
  color-gated mask).} Three attempts, documented in
  \BreakablePath{results/\discretionary{}{}{}figures/\discretionary{}{}{}mask\_\discretionary{}{}{}attempts\_\discretionary{}{}{}comparison.\discretionary{}{}{}png}
  + notebook 12 §8:

  \begin{itemize}
  \tightlist
  \item
    \textbf{(A) 2-component (bulge+disk) Sérsic + residual-only IR
    growth} --- over-masks: residual-only growth grabs red galaxy body +
    a companion (F150W2 15768 px, F322W2 4258 px) → raw biased low
    (11.16) and a large +0.31 fill-in. The bulge+disk was never physical
    for this elliptical; it was a kludge to tame the residual growth.
  \item
    \textbf{(B) single-Sérsic with a broken fit (amplitude→0) + color
    gate} --- DISCARDED artifact: the near-zero model made the residual
    term meaningless and collapsed the fill-in to \textasciitilde0.01
    mag (apparent 11.34 ± 0.013). Caught via the black model panel in
    the diagnostic figures.
  \item
    \textbf{(C) single-Sérsic (fixed init) + color-gated IR growth + WCS
    registration} --- FINAL: tighter, arc-focused, correctly-registered
    masks (F150W2 7698 px, F322W2 1997 px; companion now masked), raw
    central 11.36 \textbf{consistent with the validated expert hand-mask
    11.33}, fill real but smaller (−0.19 to −0.28), masking systematic
    ±0.086 (vs ±0.16 for 2-comp). Physically motivated (single Sérsic) +
    physically-gated (source color) + astrometrically registered
    (§2.1.1b).
  \end{itemize}
\item
  \textbf{Files} (in sibling dir
  \BreakablePath{.\discretionary{}{}{}.\discretionary{}{}{}/\discretionary{}{}{}velocity\_\discretionary{}{}{}dispersion\_\discretionary{}{}{}from\_\discretionary{}{}{}IFU/\discretionary{}{}{}}):
  \BreakablePath{AGEL020613-\discretionary{}{}{}011417A\_\discretionary{}{}{}F200LP\_\discretionary{}{}{}WFC3\_\discretionary{}{}{}cutout\_\discretionary{}{}{}L3\_\discretionary{}{}{}mask.\discretionary{}{}{}fits}
  (500×500, 25″×25″); the mask was revised 2026-02-28
  (\BreakablePath{\_\discretionary{}{}{}mask\_\discretionary{}{}{}20260228.\discretionary{}{}{}fits});
  original kept as
  \BreakablePath{\_\discretionary{}{}{}mask\_\discretionary{}{}{}OLD.\discretionary{}{}{}fits}.
\item
  \textbf{JWST bands:} same arc geometry; the HST F200LP mask is
  re-projected from the HST WCS onto each JWST frame. {[}METHODS
  §II.4{]}
\item
  \textbf{For the IFU/σ\_e application} (context): F200LP mask
  reprojected to the 0.30″ IFU grid via
  \BreakablePath{scipy.\discretionary{}{}{}ndimage.\discretionary{}{}{}map\_\discretionary{}{}{}coordinates(order=0)}
  nearest-neighbor; \textbf{35/154 = 23\% of spaxels (28\% by I-weight)}
  flagged inside R\_e=2.097″ (a geometric LOWER bound --- the 1.27″
  seeing spreads the arc much wider; see §2.4.1 step 3 +
  \BreakablePath{scripts/\discretionary{}{}{}psf\_\discretionary{}{}{}mask\_\discretionary{}{}{}fraction.\discretionary{}{}{}py}).
  {[}feedback\_masking\_strategy.md{]}
\end{itemize}

\hypertarget{estimating-r_e}{%
\subsubsection{§2.1.2 Estimating R\_e}\label{estimating-r_e}}

\begin{itemize}
\tightlist
\item
  \textbf{Headline R\_e = 2.097″ = 14.76 kpc} at z = 0.67564 (best-mask
  F140W+F200LP CoG). {[}CLAUDE.md; METHODS §II.3/A3{]}
\item
  \textbf{``Best mask throughout'' (2026-06-11):} the headline R\_e is
  now the \textbf{best-mask} (single-Sérsic + color/morph gate + WCS
  registration) curve-of-growth = \textbf{2.097″}, superseding the
  expert-mask 2.305″. The best masks genuinely remove outer diffuse-arc
  + interloper light that slightly inflated the expert-mask R\_e; the
  model-based Sérsic r\_eff (1.897″) and sky-subtracted CoG (1.922″)
  both confirm the inward shift. The HST 2-R\_e companion mask is empty
  (companions live only in deep JWST), so it does not enter the HST R\_e
  --- the global color/morph mask alone gives 2.097″ (self-consistent
  fixed point).
  \BreakablePath{scripts/\discretionary{}{}{}re\_\discretionary{}{}{}mask\_\discretionary{}{}{}sensitivity.\discretionary{}{}{}py},
  \BreakablePath{scripts/\discretionary{}{}{}Re\_\discretionary{}{}{}bestmask\_\discretionary{}{}{}reconciliation.\discretionary{}{}{}py}.
\item
  \textbf{photutils validation
  (\BreakablePath{scripts/\discretionary{}{}{}validate\_\discretionary{}{}{}Re\_\discretionary{}{}{}photutils.\discretionary{}{}{}py}):}
  our azimuthally-averaged masked CoG is independently reproduced by
  photutils \BreakablePath{RadialProfile} (integrated) to
  \textbf{±0.002″}, and unmasked matches photutils
  \BreakablePath{CurveOfGrowth} (direct 2D aperture sum) to ±0.004″. A
  naive masked aperture-\emph{sum} (photutils
  \BreakablePath{CurveOfGrowth} dropping masked pixels) would bias R\_e
  \textbf{+0.25--0.41″ high} → our azimuthal-fill is the correct
  masked-CoG treatment.
  \BreakablePath{centroid\_\discretionary{}{}{}2dg} is robust;
  \BreakablePath{centroid\_\discretionary{}{}{}quadratic} fails on
  extended light (expected --- a point-source tool).
\item
  \textbf{Method systematic (best mask) = ±0.100″} (raw CoG 2.097 /
  sky-sub CoG 1.922 / Sérsic r\_eff 1.897, peak-to-peak/2;
  \BreakablePath{results/\discretionary{}{}{}Re\_\discretionary{}{}{}cog\_\discretionary{}{}{}reconciliation\_\discretionary{}{}{}bestmask.\discretionary{}{}{}npz}).
\item
  \textbf{Method:} mean of the \textbf{F140W and F200LP best-mask
  curves-of-growth}
  (\BreakablePath{scripts/\discretionary{}{}{}final\_\discretionary{}{}{}sigma\_\discretionary{}{}{}e.\discretionary{}{}{}py:curve\_\discretionary{}{}{}of\_\discretionary{}{}{}growth},
  the PRODUCTION script):

  \begin{itemize}
  \tightlist
  \item
    best-mask F140W CoG = \textbf{1.912″}; F200LP CoG = \textbf{2.281″};
    mean = \textbf{2.097″} (expert-mask values were 2.168/2.441/2.305).
  \item
    Masked CoG zeros mask-flagged HST pixels before the radial flux
    integral, so the lensed arc / companions do not bias the half-light
    point.
  \item
    Centers: HST-mean
    \BreakablePath{photutils.\discretionary{}{}{}centroid\_\discretionary{}{}{}2dg}
    centers (same as the σ\_e pipeline), both bands.
  \end{itemize}
\item
  \textbf{Do NOT confuse with
  \BreakablePath{scripts/\discretionary{}{}{}measure\_\discretionary{}{}{}Re.\discretionary{}{}{}py}}
  →
  \BreakablePath{results/\discretionary{}{}{}Re\_\discretionary{}{}{}measurements.\discretionary{}{}{}npz}.
  That script uses the image-geometric center and produces \textbf{10
  variants} (3 sources × 4 masking strategies: \BreakablePath{unmasked},
  \BreakablePath{zeroed}, \BreakablePath{proper},
  \BreakablePath{PSF-\discretionary{}{}{}convolved}); its
  ``proper-mask'' mean (2.633″) is \textbf{not the headline} ---
  early-reference only.
  (\BreakablePath{measure\_\discretionary{}{}{}Re.\discretionary{}{}{}hst\_\discretionary{}{}{}Re}
  previously hardcoded pixscale=0.08 for both bands; \textbf{fixed
  2026-06-08 (A3c/nb14) to read the WCS pixscale +
  \BreakablePath{r\_\discretionary{}{}{}max\_\discretionary{}{}{}arcsec=6.\discretionary{}{}{}0},
  now reconciling with the production CoG to \textless0.04″} --- but
  \BreakablePath{final\_\discretionary{}{}{}sigma\_\discretionary{}{}{}e}
  remains the headline path.)
\item
  \textbf{Superseded prior values:} IFU white-light R\_e = 2.61″
  (nb05/06 era); expert-mask HST CoG = 2.305″ (pre-2026-06-11). A Ca
  H+K+G depth-map gives 2.90″ (kinematic-tracer estimator, not the
  photometric R\_e).
\item
  \textbf{R\_e as a σ\_e systematic (test D7 → best-mask grid):} σ\_e
  rises monotonically with R\_e. Re-measured at the wide window over a
  7-point R\_e grid bracketing the 2.097″ headline
  (\BreakablePath{scripts/\discretionary{}{}{}run\_\discretionary{}{}{}sigma\_\discretionary{}{}{}e\_\discretionary{}{}{}Re\_\discretionary{}{}{}grid.\discretionary{}{}{}py});
  the \textbf{adopted} R\_e-source systematic = \textbf{±5.13 km/s}
  (best-mask CoG family \{1.912/2.097/2.281″\}), folded into the budget
  (§3.1, §2.4.3). The full grid incl.~CaHK+G 2.90″ (±9.98) and CaHK+G's
  +12.4 km/s deviation are cross-checks, NOT folded. {[}CLAUDE.md{]}
\end{itemize}

\hypertarget{measuring-fluxes}{%
\subsubsection{§2.1.3 Measuring Fluxes}\label{measuring-fluxes}}

\begin{itemize}
\tightlist
\item
  \textbf{Tools:} \BreakablePath{photutils} aperture photometry
  (interactive masking GUIs
  \BreakablePath{scripts/\discretionary{}{}{}photometry\_\discretionary{}{}{}masking\_\discretionary{}{}{}\{HST,JWST\}.\discretionary{}{}{}py}).
  \textbf{\BreakablePath{synphot} is NOT currently used} --- flux→AB is
  done directly from header zeropoints. \emph{(If the outline wants
  synphot, that is aspirational; flag with co-authors.)} {[}METHODS
  §II.4{]}
\item
  \textbf{AB zeropoints --- read per-image from FITS headers, never
  hardcoded.} {[}CLAUDE.md; reference\_hst\_jwst\_data\_properties.md{]}

  \begin{itemize}
  \tightlist
  \item
    HST (drizzled e⁻/s):
    \BreakablePath{ZP\_\discretionary{}{}{}AB = −2.\discretionary{}{}{}5·log10(PHOTFLAM) − 21.\discretionary{}{}{}10 − 5·log10(PHOTPLAM) + 18.\discretionary{}{}{}6921}.
  \item
    JWST (i2d, MJy/sr):
    \BreakablePath{ZP\_\discretionary{}{}{}AB = −6.\discretionary{}{}{}10 − 2.\discretionary{}{}{}5·log10(PIXAR\_\discretionary{}{}{}SR)}.
  \end{itemize}
\item
  \textbf{Four bands / instruments / programs:}
\end{itemize}

\begin{longtable}[]{@{}
  >{\raggedright\arraybackslash}p{(\columnwidth - 14\tabcolsep) * \real{0.12}}
  >{\raggedright\arraybackslash}p{(\columnwidth - 14\tabcolsep) * \real{0.12}}
  >{\raggedright\arraybackslash}p{(\columnwidth - 14\tabcolsep) * \real{0.12}}
  >{\raggedright\arraybackslash}p{(\columnwidth - 14\tabcolsep) * \real{0.12}}
  >{\raggedright\arraybackslash}p{(\columnwidth - 14\tabcolsep) * \real{0.12}}
  >{\raggedright\arraybackslash}p{(\columnwidth - 14\tabcolsep) * \real{0.12}}
  >{\raggedright\arraybackslash}p{(\columnwidth - 14\tabcolsep) * \real{0.12}}
  >{\raggedright\arraybackslash}p{(\columnwidth - 14\tabcolsep) * \real{0.12}}@{}}
\toprule
Telescope & Instrument & Filter & Program (PI) & Drizzled scale &
EXPTIME & ZP\_AB & AB (aperture) \\
\midrule
\endhead
HST & WFC3/UVIS & F200LP & 16773 (Glazebrook) & 0.0500″/pix & 600.0 s &
27.344 & 22.613 \\
HST & WFC3/IR & F140W & 16773 (Glazebrook) & 0.0800″/pix & 597.7 s &
26.446 & 19.134 \\
JWST & NIRCam SW & F150W2 & 05594 (Mahler) & 0.03075″/pix & 1836.0 s &
28.033 & 18.942 \\
JWST & NIRCam LW & F322W2 & 05594 (Mahler) & 0.06301″/pix & 1836.0 s &
26.475 & 18.604 \\
\bottomrule
\end{longtable}

{[}reference\_hst\_jwst\_data\_properties.md{]} - JWST pointing comes
from the cluster program ACT-CL J0206.2−0114 (ACT DR5, Hilton et
al.~2021); the deflector sits in that field. No PAM correction
(drizzled). \textbf{Drizzle pixel correlation:} empirical noise
\textasciitilde1.3--2× the idealized CCD-equation value --- motivates
the 10\% flux floor (§2.1.4). - \textbf{2D Sérsic verification given
masking (the morphology cross-check; notebook 08):} single-component 2D
Sérsic per band
(\BreakablePath{scripts/\discretionary{}{}{}measure\_\discretionary{}{}{}Re.\discretionary{}{}{}py}
framework + the \textbf{2026-05-11 bound-fix} n∈{[}1.0,6.0{]},
ellip∈{[}0.0,0.6{]}, 3-init grid --- prevented an n=0.30 flat-disk
escape on F200LP), extrapolated to total flux:

\begin{longtable}[]{@{}llllll@{}}
\toprule
Filter & AB aperture & AB Sérsic-total & Δmag & n & r\_eff (″) \\
\midrule
\endhead
F200LP & 22.613 & 20.672 & −1.94 & 1.42 & 2.49 \\
F140W & 19.134 & 18.282 & −0.85 & 1.54 & 1.88 \\
F150W2 & 18.942 & 18.154 & −0.79 & 1.40 & 1.72 \\
F322W2 & 18.604 & 17.633 & −0.97 & 1.97 & 2.03 \\
\bottomrule
\end{longtable}

Aperture under-counts outer-profile flux by 0.8--1.9 mag, but the
integrated-mass effect is only +0.07 dex (see §2.1.4 / Results). Caches:
\BreakablePath{results/\discretionary{}{}{}sersic\_\discretionary{}{}{}total\_\discretionary{}{}{}photometry.\discretionary{}{}{}npz},
\BreakablePath{results/\discretionary{}{}{}bagpipes\_\discretionary{}{}{}sersic\_\discretionary{}{}{}refit.\discretionary{}{}{}npz}.
{[}METHODS §II.6{]}

\hypertarget{sed-modeling-with-bagpipes}{%
\subsubsection{\texorpdfstring{§2.1.4 SED Modeling with
\BreakablePath{bagpipes}}{§2.1.4 SED Modeling with }}\label{sed-modeling-with-bagpipes}}

\begin{itemize}
\item
  \textbf{Tool:} \BreakablePath{bagpipes} (Carnall et al.~2018, 2019);
  BPASS+CLOUDY templates (Bagpipes default). Notebook
  \BreakablePath{02\_\discretionary{}{}{}streamlined\_\discretionary{}{}{}Bagpipes\_\discretionary{}{}{}SED.\discretionary{}{}{}ipynb}.
  {[}METHODS §II.5; reference\_sed\_fitting.md{]}
\item
  \textbf{Model:} exponential-τ SFH; Calzetti dust (A\_V free, 0--2);
  free metallicity (0--2.5 Z☉); redshift prior (0.674, 0.676) bracketing
  the line-fit z. Sampler \textbf{Nautilus} (n\_live=400), 500-sample
  posterior.

\begin{Shaded}
\begin{Highlighting}[]
\NormalTok{fit\_instructions }\OperatorTok{=}\NormalTok{ \{}
  \StringTok{"redshift"}\NormalTok{:   (}\FloatTok{0.674}\NormalTok{, }\FloatTok{0.676}\NormalTok{),}
  \StringTok{"exponential"}\NormalTok{: \{}\StringTok{"age"}\NormalTok{: (}\FloatTok{0.1}\NormalTok{,}\FloatTok{15.}\NormalTok{), }\StringTok{"tau"}\NormalTok{: (}\FloatTok{0.3}\NormalTok{,}\FloatTok{10.}\NormalTok{),}
                  \StringTok{"massformed"}\NormalTok{: (}\FloatTok{1.}\NormalTok{,}\FloatTok{15.}\NormalTok{), }\StringTok{"metallicity"}\NormalTok{: (}\FloatTok{0.}\NormalTok{,}\FloatTok{2.5}\NormalTok{)\},}
  \StringTok{"dust"}\NormalTok{: \{}\StringTok{"type"}\NormalTok{: }\StringTok{"Calzetti"}\NormalTok{, }\StringTok{"Av"}\NormalTok{: (}\FloatTok{0.}\NormalTok{, }\FloatTok{2.}\NormalTok{)\},}
\NormalTok{\}}
\end{Highlighting}
\end{Shaded}
\item
  \textbf{Why a 10\% fractional error floor on each band:} a
  conservative, uniform systematic budget that covers both the formal
  photon noise \textbf{and} the drizzle pixel-correlation underestimate
  (empirical noise \textasciitilde1.3--2× the idealized CCD-equation
  value). {[}METHODS §II.4/II.7; reference\_sed\_fitting.md{]}
\item
  \textbf{Systematic test to fill in masked pixels = the 2D-Sérsic-total
  refit:} because the aperture excludes arc-masked pixels and truncates
  the outer profile, we fit a 2D Sérsic per band (§2.1.3), extrapolate
  to total flux, and re-run Bagpipes with identical priors. This
  recovers the flux missing from masked/truncated regions and bounds the
  mass bias. Result: log(M⋆/M☉) = \textbf{11.40 +0.11/−0.15}
  (Sérsic-total) vs \textbf{11.33 +0.07/−0.09} (aperture) → aperture is
  biased low by \textasciitilde0.07 dex (+16\% in M⋆), within the quoted
  uncertainty. {[}METHODS §II.6{]} \textbf{NOTE (updated 2026-06-10):}
  this analytic-total refit is superseded by the per-aperture
  \textbf{single-Sérsic fill-in} under the principled color-gated
  IR-extended masks (§2.1.1b / §3.2): the headline is the empirical raw
  value \textbf{11.36 (10\%)} with a one-sided \textbf{+0.14} dex
  fill-in systematic reaching \textbf{11.47}. The 11.40 analytic-total
  sits just inside that bracket and the central is consistent with the
  validated expert hand-mask 11.33. The enhanced comparison figure
  \BreakablePath{results/\discretionary{}{}{}figures/\discretionary{}{}{}nb08\_\discretionary{}{}{}sersic\_\discretionary{}{}{}vs\_\discretionary{}{}{}aperture.\discretionary{}{}{}png}
  now overlays, for both the empirical-aperture and Sérsic-total
  choices, the best-fit Bagpipes model SED, the measured photometry, and
  the filter-convolved model photometry.
\item
  \textbf{Posteriors (older expert-aperture, now a cross-check):}
  log(M⋆/M☉)=11.33 +0.07/−0.09; mass-weighted age 5.10 Gyr
  {[}3.69--6.11{]}; τ = 0.68 Gyr; Z = 0.89 Z☉; A\_V = 0.82 mag; z\_phot
  = 0.675 {[}0.674--0.676{]}. Per-band residuals: F200LP +4.7\%, F140W
  −11.3\%, F150W2 −8.0\%, F322W2 +9.6\% (mild single-τ-SFH NIR-slope
  tension; not catastrophic, within ±0.07--0.15 dex). {[}METHODS
  §II.5{]}
\item
  \textbf{Caches / figure:}
  \BreakablePath{results/\discretionary{}{}{}bagpipes\_\discretionary{}{}{}sed\_\discretionary{}{}{}results.\discretionary{}{}{}npz}
  (the file Figure 2 loads),
  \BreakablePath{pipes/\discretionary{}{}{}posterior/\discretionary{}{}{}AGEL0206/\discretionary{}{}{}0206\_\discretionary{}{}{}real.\discretionary{}{}{}h5}
  (delete to refit),
  \BreakablePath{results/\discretionary{}{}{}bagpipes\_\discretionary{}{}{}sersic\_\discretionary{}{}{}refit.\discretionary{}{}{}npz}.
\end{itemize}

\begin{center}\rule{0.5\linewidth}{0.5pt}\end{center}

\hypertarget{spectroscopic-analysis}{%
\subsection{§2.2 Spectroscopic Analysis}\label{spectroscopic-analysis}}

\hypertarget{calibration-keerthi-kaustubh}{%
\subsubsection{§2.2.1 Calibration --- @Keerthi /
Kaustubh}\label{calibration-keerthi-kaustubh}}

\begin{itemize}
\tightlist
\item
  \textbf{Instrument:} Keck II / KCWI, \textbf{KCRM red arm}, Red-Low
  (RL) grating, Medium slicer, 2×2 binning; dithered at PA 0°/45°/90°.
  {[}METHODS §I.1; reference\_kcwi\_data\_properties.md{]}
\item
  \textbf{Cube (headline):}
  \BreakablePath{Nov17\_\discretionary{}{}{}2025\_\discretionary{}{}{}DESJ0206\_\discretionary{}{}{}RL\_\discretionary{}{}{}combined\_\discretionary{}{}{}mtwdo\_\discretionary{}{}{}icubes\_\discretionary{}{}{}wcs.\discretionary{}{}{}fits}
  --- the NEW
  \textbf{\BreakablePath{\_\discretionary{}{}{}mtwdo\_\discretionary{}{}{}}}
  reduction (Master Twilight + Dome hybrid flats), K. R. Gupta's latest
  re-reduction. The earlier
  no-\BreakablePath{\_\discretionary{}{}{}mtwdo\_\discretionary{}{}{}}
  cube is the OLD reduction (second-reduction cross-check; co-aligned to
  \textless0.002″, differs only in flat-fielding; σ\_e differs
  \textasciitilde7 km/s --- the ``reduction-pass'' systematic).
\item
  \textbf{Reduction chain:} raw frames → \BreakablePath{kcwidrp};
  stacked with
  \BreakablePath{ISMgas/\discretionary{}{}{}kcwi/\discretionary{}{}{}kcwiReduction}
  by \textbf{K. R. Gupta (Kaustubh)}; Dec-29 RED frames pre-reduced by
  Yuguang Chen. \emph{(Calibration prose is Keerthi/Kaustubh's to write
  --- these are the facts on hand.)}
\item
  \textbf{Cube geometry:} shape (3317, 100, 100); \textbf{0.300″
  spaxels}; \textasciitilde30″×30″ FoV; λ = 5625--8941 Å at Δλ = 1.0
  Å/pix; BUNIT = FLAM16/arcsec².
\item
  \textbf{Three confirmed nights} (totals ≈ \textbf{180 min RED + 198
  min BLUE} on source):
\end{itemize}

\begin{longtable}[]{@{}llll@{}}
\toprule
Night (UT) & Program (PI) & RED & BLUE \\
\midrule
\endhead
2024 Dec 29 & U204 & 4×300 s & 1×1320 s (excluded by reducer) \\
2025 Aug 30 & K409 (PI TBD) & 12×300 s & 4×990 s \\
2025 Nov 17 & U002 (Jones) & 20×300 s & 5×1320 s \\
\bottomrule
\end{longtable}

\textbf{FLAG:} the cube header (DATE-OBS 2025-11-17, PROGPI Jones)
describes only the first input frame --- cite the multi-night combine,
not the header. K409 PI and Aug-30 DIMM seeing still TBD. -
\textbf{Spectral resolution / LSF:} FWHM\_inst = 2.355 ×
DISPSCAL(0.2940) × 1.0 Å = \textbf{0.692 Å} → \textbf{R ≈ 10000} at the
fit band; \textbf{σ\_v,inst ≈ 12.6 km/s observed (7.5 km/s rest-frame)}
--- the \textasciitilde270 km/s dispersion is resolved by
\textasciitilde20×. LSF robustness ≤ 0.83 km/s over DISPSCAL ×0.5--×2.0.
- \textbf{Seeing:} guide-camera FWHM = \textbf{1.27″}
(\BreakablePath{GUIDFWHM}, Nov-17; conservative).

\hypertarget{redshift-measurement}{%
\subsubsection{§2.2.2 Redshift Measurement}\label{redshift-measurement}}

\begin{itemize}
\tightlist
\item
  \textbf{Method --- single/multi-line Gaussian fits} to the
  integrated-spectrum features on \textbf{NIST air rest wavelengths},
  per-line z inverse-variance combined
  (\BreakablePath{scripts/\discretionary{}{}{}redshift\_\discretionary{}{}{}verify.\discretionary{}{}{}py},
  notebook 04).
\end{itemize}

\textbf{Deflector absorption features (systemic z = 0.67564 ± 0.00033; 6
fits):}

\begin{longtable}[]{@{}llllll@{}}
\toprule
Feature & Ion/species & rest λ\_air (Å) & obs λ\_air (Å) & z\_line & σ
(Å) \\
\midrule
\endhead
Ca K & Ca II & 3933.66 & 6591.79 & 0.67574 ± 0.00096 & 11.5 \\
Ca H & Ca II & 3968.47 & 6649.61 & 0.67561 ± 0.00097 & 16.6 \\
Hδ & H I & 4101.74 & 6871.80 & 0.67534 ± 0.00058 & 3.3 \\
G-band & CH (band) & 4304.40 & 7209.31 & 0.67487 ± 0.00160 & 11.3 \\
Hβ & H I & 4861.33 & 8149.30 & 0.67635 ± 0.00082 & 2.8 \\
Ca H+K & Ca II doublet & 3933.66 / 3968.47 & --- & 0.67567 ± 0.00069 &
11.2 \\
\textbf{Weighted} & & & & \textbf{0.67564 ± 0.00033} & \\
\bottomrule
\end{longtable}

\textbf{Source emission features (lensed arc, z = 1.30263 ± 0.00003;
{[}O II{]} doublet, red cube):}

\begin{longtable}[]{@{}llll@{}}
\toprule
Feature & rest λ\_air (Å) & obs λ\_air (Å) & z\_line \\
\midrule
\endhead
{[}O II{]} 3726 & 3726.03 & 8579.2 & 1.30357 ± 0.00005 \\
{[}O II{]} 3729 & 3728.82 & 8585.6 & 1.30185 ± 0.00005 \\
\textbf{{[}O II{]} doublet} & & & \textbf{1.30251 ± 0.00004} \\
\bottomrule
\end{longtable}

\begin{itemize}
\tightlist
\item
  \textbf{Deflector systemic z = 0.67564} (supersedes AGEL DR2 z =
  0.67511; \textasciitilde95 km/s offset, used only in old notebooks).
  {[}METHODS §I.2; reference\_redshift\_verification.md{]}
\item
  Other absorption features in
  \BreakablePath{ABSORPTION\_\discretionary{}{}{}LINES} (Hγ 4340.47,
  Fe4383 4383.55, Mg I/Mg b 5175.27, Fe I 5270.40, Na D 5893.00) are
  NIST-air but were NOT used in the systemic fit (out of the red-cube
  range at z=0.676, or blended); listed for completeness in
  \BreakablePath{scripts/\discretionary{}{}{}redshift\_\discretionary{}{}{}verify.\discretionary{}{}{}py}.
\end{itemize}

\begin{quote}
\textbf{Air↔vacuum audit (2026-06-11 --- PASS).} (1) Every rest
wavelength is the \textbf{NIST air} value, verified to
\textbf{\textless0.01 Å} (Ca K 3933.66 vs NIST 3933.663, Hβ 4861.33 vs
4861.333, {[}O II{]} 3726.03 vs 3726.032, etc.). \textbf{G-band 4304.40
is the CH molecular bandhead} (a blend centroid, not a single atomic
line) → it is the most discrepant/uncertain line (z=0.67487, ±0.0016)
and is down-weighted by its error. (2) The KCWI cube is \textbf{vacuum}
(FITS \BreakablePath{CTYPE3='WAVE'}; \BreakablePath{AWAV} would denote
air) and is converted to \textbf{air} before fitting via the
\textbf{Ciddor (1996)/Morton (2000) IAU-standard} index of refraction n
= 1 + 8.34254×10⁻⁵ + 2.406147×10⁻²/(130−s²) + 1.5998×10⁻⁴/(38.9−s²), s =
10⁴/λ\_vac --- \textbf{the same formula the σ\_e ppxf pipeline uses}
(\BreakablePath{bootstrap\_\discretionary{}{}{}ppxf.\discretionary{}{}{}py};
Ciddor 1996). (3) Thus z = λ\_obs,air/λ\_rest,air − 1 combines
\textbf{both wavelengths in the air frame} → internally consistent, z
\textbf{unbiased}. Had we left the cube in vacuum and used air rest
lines, z would be biased high by ≈ (1+z)·(n−1) ≈ +5×10⁻⁴ (≈ +90 km/s);
the conversion removes this. - \textbf{Double verification with ppxf:}
the Gaussian line-fit z is independent of, and consistent with, the ppxf
V\_sys (ppxf returns z via z = (1+z₀)·exp(V/c) − 1; Cappellari 2023, the
ppxf relativistic velocity convention --- \emph{eq. number not yet
verified, do not assert ``eq. 5c''}). - \textbf{Source z = 1.30263 ±
0.00003} --- measured from the \textbf{{[}O II{]} λλ3726,3729 doublet}
in the red cube (notebook 04; obs 8579.2/8585.6 Å), consistent with AGEL
DR2 z = 1.302. (The ±0.00003 is the formal doublet-fit error; the
per-line spread 1.30185--1.30357 is the realistic uncertainty.) -
\textbf{Source rest-UV resonance lines (arc ISM;
\BreakablePath{scripts/\discretionary{}{}{}source\_\discretionary{}{}{}uv\_\discretionary{}{}{}redshift.\discretionary{}{}{}py},
2026-06-11).} The bright star-forming arc shows the classic Mg II / Fe
II resonance absorption: - \textbf{Mg II λλ2796.35,2803.53 (red arm):}
strongly detected (\textbf{16.7σ / 18.4σ}), z(absorption) =
\textbf{1.30113 ± 0.00003}, i.e.~a \textbf{−195 km/s} centroid offset vs
the systemic {[}O II{]} z=1.30263 (both in the \emph{same} red cube, so
it is a real \emph{relative} offset, not cross-arm calibration).
\textbf{We do NOT interpret this as an outflow} without further work ---
for a resonance line on a \emph{lensed} arc the offset is degenerate
among: (i) Mg II \textbf{emission infill} filling the red absorption
wing (P-Cygni; pulls the centroid blueward with no bulk motion), (ii)
\textbf{differential lensing / aperture} --- the Mg II absorption
(continuum-weighted) and {[}O II{]} emission (H II-region-weighted)
sample different source regions that the lens positions/magnifies
differently, and (iii) the blended-doublet single-Gaussian centroid on
an asymmetric profile. Report it as an \textbf{observed Mg II--{[}O
II{]} velocity offset}, with outflow as one (unconfirmed) possibility. -
\textbf{Fe II λλ2344.21,2382.77 (blue arm):} detected (7.9σ / 3.0σ; Fe
II 2374 only 1.2σ), but at z ≈ 1.306, \textasciitilde+470 km/s from the
red systemic. Because Fe II sits in the \textbf{separate blue cube} (an
older \BreakablePath{AWAV}/air reduction, ≠ the red
\BreakablePath{WAVE}/vacuum mtwdo cube), this offset is dominated by a
\textbf{blue↔red wavelength-calibration difference}, so the Fe II
\emph{velocity} is unreliable cross-arm --- the detection is real, the
velocity is a
\BreakablePath{[TODO: cross-\discretionary{}{}{}calibrate the blue arm]}.
(Rest λ are Morton 2003 vacuum; the blue cube is converted air→vac for
the fit. NOT a clean independent z --- quote Mg II / {[}O II{]}.) -
\textbf{Companion ellipticals --- a galaxy GROUP at the deflector
redshift} (notebook 04a,
\BreakablePath{scripts/\discretionary{}{}{}redshift\_\discretionary{}{}{}verify.\discretionary{}{}{}bootstrap\_\discretionary{}{}{}redshift},
N=200 wild bootstrap on the red cube). Two early-type companions flank
the primary deflector, \textbf{both at z ≈ 0.676 = the deflector z}
(physically associated, not interlopers):
\end{quote}

\begin{longtable}[]{@{}
  >{\raggedright\arraybackslash}p{(\columnwidth - 8\tabcolsep) * \real{0.20}}
  >{\raggedright\arraybackslash}p{(\columnwidth - 8\tabcolsep) * \real{0.20}}
  >{\raggedright\arraybackslash}p{(\columnwidth - 8\tabcolsep) * \real{0.20}}
  >{\raggedright\arraybackslash}p{(\columnwidth - 8\tabcolsep) * \real{0.20}}
  >{\raggedright\arraybackslash}p{(\columnwidth - 8\tabcolsep) * \real{0.20}}@{}}
\toprule
Companion & position & sep. from deflector & z (absorption) & z
(emission) \\
\midrule
\endhead
\textbf{NE} (``2 o'clock'', larger) & IFU (x,y)=(61,64) & ≈ 4.5″ ≈ 31
kpc NE & 0.67580 −0.00067/+0.00062 & 0.67565 −0.00045/+0.00060 \\
\textbf{SW} (``7 o'clock'', compact) & IFU (x,y)=(38,45) & ≈ 4.6″ ≈ 32
kpc SW & 0.67588 −0.00034/+0.00047 & 0.67558 −0.00048/+0.00071 \\
\bottomrule
\end{longtable}

Both companions sit within \textbar Δcz\textbar{} ≲ 60 km/s of the
deflector (0.67564) → AGEL J020613−011417 is the central galaxy of a
\textbf{compact group/triple} at z ≈ 0.676. (These are the field
ellipticals masked in the aperture photometry; their group membership
confirms they are not background interlopers.) - \textbf{Peculiar
velocities:} we \textbf{do not} correct for peculiar velocities --- we
use the line-fit systemic z directly (see gap G3; this sentence needs to
be added to the draft). - \textbf{Wavelength convention:} AIR rest
wavelengths throughout (NIST-verified to \textless0.01 Å); KCWI native
vacuum (\BreakablePath{CTYPE3='WAVE'}) converted to air via Ciddor 1996
--- see the air↔vacuum audit box above. Do not mix SDSS vacuum values
(\textasciitilde1 Å ≈ 80 km/s systematic at 4000 Å).
{[}feedback\_wavelength\_convention.md{]}

\hypertarget{stellar-kinematics-velocity-dispersion}{%
\subsubsection{§2.2.3 Stellar Kinematics \& Velocity
Dispersion}\label{stellar-kinematics-velocity-dispersion}}

\textbf{Spatial regime / masking / R\_e}

\begin{itemize}
\tightlist
\item
  Headline σ\_e is measured on a \textbf{single I(r)-weighted 1-D
  aperture spectrum at R \textless{} R\_e} --- the co-add-then-fit
  method of \textbf{Cappellari et al.~2006 (SAURON IV, §2.3)}, the
  literature-standard σ\_e (KH13 / Greene+20 / SAURON / ATLAS3D / MaNGA
  convention). \textbf{R\_e = 2.097″ = 14.76 kpc} (best mask; §2.1.2).
\item
  \textbf{Center:} HST-mean
  \BreakablePath{photutils.\discretionary{}{}{}centroid\_\discretionary{}{}{}2dg}
  of F140W+F200LP cores, propagated through the KCWI WCS; adopted RA =
  31.55613°, Dec = −1.23819° (ICRS). F140W↔F200LP offset 0.36″ (F200LP
  shifted by the UV-bright arc; F140W is the clean bulge centroid).
\item
  \textbf{Spatial arc mask:} F200LP segmentation mask reprojected to the
  IFU grid (nearest-neighbor); 38/184 spaxels inside R\_e flagged
  (\textasciitilde27\% by I-weight). F140W mask not used (covers the
  core).
\item
  \textbf{I-weight map:} per-spaxel collapse of the IFU 6500--7500 Å
  white-light band; spectra \textbf{summed} (not averaged) so the wild
  bootstrap preserves per-spaxel noise.
\item
  R \textless{} R\_e/8 (JFK95 σ\_c) was \textbf{dropped} --- at 0.288″
  it sits inside seeing FWHM/2 = 0.64″ (\textasciitilde3 spaxels). No
  central σ is quoted; σ\_e(\textless R\_e/2) ≈ 225 km/s is reported as
  a gradient diagnostic only.
\end{itemize}

\textbf{ppxf setup}

\begin{itemize}
\tightlist
\item
  \textbf{What ppxf does} --- \textbf{Penalized Pixel-Fitting}
  (Cappellari \& Emsellem 2004, PASP 116, 138; upgraded in Cappellari
  2017, MNRAS 466, 798; Cappellari 2023 --- \emph{volume/page not
  re-verified here}). It forward-models the observed galaxy spectrum, in
  \textbf{pixel space on a log-rebinned (constant km/s per pixel) grid},
  as a \textbf{non-negative linear combination of stellar templates
  convolved with a parametrised line-of-sight velocity distribution
  (LOSVD)}, plus low-order additive Legendre polynomials that absorb
  continuum/template mismatch. The template weights are solved by
  non-negative least-squares and the LOSVD parameters by
  Levenberg--Marquardt, alternately, to minimise χ². The
  \emph{``penalized''} is an optional bias that pulls the higher
  Gauss--Hermite moments toward a Gaussian at low S/N ---
  \textbf{inactive here}, since we fit a pure Gaussian LOSVD.
\item
  \textbf{ppxf v9.4.5}, \textbf{\BreakablePath{moments = 2}} (fit only V
  and σ; no h₃/h₄), \BreakablePath{bias} left at the default (irrelevant
  at moments=2), \BreakablePath{clean = False} (cosmic rays are handled
  by the
  \BreakablePath{BAD\_\discretionary{}{}{}PIXELS\_\discretionary{}{}{}REST}
  mask instead --- ppxf's internal σ-clip rejects 0 pixels here because
  the KCWI noise array is over-estimated vs the residual scatter),
  \BreakablePath{velscale} set to the cube's native sampling,
  \BreakablePath{trig = False}. {[}scripts/bootstrap\_ppxf.py,
  \BreakablePath{run\_\discretionary{}{}{}wide\_\discretionary{}{}{}sigma\_\discretionary{}{}{}e.\discretionary{}{}{}py};
  METHODS §I.5/I.6{]}
\item
  \textbf{Polynomial degrees:} additive Legendre \BreakablePath{degree}
  swept over \textbf{15--29 (15 values)} and pooled;
  \BreakablePath{mdegree = 0}. σ-vs-degree is flat within the bootstrap
  envelope (no LOSVD absorption by the polynomial). \emph{(Note: METHODS
  §I.5 prose mis-labels these ``multiplicative'' --- the code passes
  \BreakablePath{degree=…, mdegree=0}, i.e.~additive. Describe them as
  additive.)}
\item
  \textbf{Three SPS libraries, pooled at the posterior level:} FSPS,
  EMILES, XSL (8--9 templates each over rest 3500--5000 Å).
\item
  \textbf{Per-SPS native wavelength frame:} FSPS = vacuum, EMILES = air,
  XSL = air. Galaxy frame matched per-SPS (scalar-median vac↔air via
  Ciddor 1996, following Cappellari's pattern). Why: the legacy
  ``convert galaxy to air for all three'' produced a spurious −90 km/s
  FSPS V\_sys; the fix collapsed the V\_sys split-track
  \textasciitilde110 → \textasciitilde15--18 km/s (σ shift ≤2 km/s).
  {[}reference\_ppxf\_vacair\_handling.md{]}
\item
  \textbf{Per-SPS V\_sys subtracted before pooling} (σ is V-invariant; a
  single global V\_sys would inflate σ\_e by \textasciitilde10--15 km/s
  in the discrete cross-check). \textbf{All three libraries now converge
  to the same redshift} --- at the headline R\_e=2.097″ aperture the
  per-SPS V\_sys are FSPS −15.7 / EMILES −19.1 / XSL −14.9 km/s → z =
  0.67555 / 0.67553 / 0.67556 (a \textasciitilde4 km/s, Δz≈3×10⁻⁵
  spread; mean \textbf{z\_ppxf = 0.67555 ± 0.00001}). This convergence
  is \emph{the} signature that the air↔vacuum frame fix worked: the
  legacy ``convert all three to air'' left FSPS (the vacuum-native
  library) \textasciitilde90--110 km/s discrepant from EMILES/XSL;
  per-SPS native framing collapses it to the \textasciitilde4 km/s seen
  here (residual folded into the ±5 km/s frame systematic).
\item
  \textbf{Emission-line handling (sigma clipping / removing emission
  lines):}

  \begin{itemize}
  \tightlist
  \item
    Custom \textbf{no-Balmer goodpixels}
    (\BreakablePath{\_\discretionary{}{}{}determine\_\discretionary{}{}{}goodpixels\_\discretionary{}{}{}no\_\discretionary{}{}{}balmer}):
    masks only the forbidden lines ({[}O II{]}, {[}O III{]}, {[}O I{]},
    {[}N II{]}, {[}S II{]}) and \textbf{keeps Hδ, Hγ, Hβ IN the fit}
    (they are absorption in this passive deflector; masking them
    discards \textasciitilde120 stellar pixels). \textbf{Hδ: keep
    unmasked --- RESOLVED (M12/nb16):} Hδ is well-fit (local-MAD not an
    outlier), and masking it shifts σ\_e by +6--8 km/s = LOSVD
    \emph{information}, not contamination (the M9 ``it's signal''
    pattern); the
    \BreakablePath{bootstrap\_\discretionary{}{}{}ppxf.\discretionary{}{}{}py}
    TODO is closed.
  \item
    \textbf{z = 1.302 source-emission masks}
    (\BreakablePath{ARC\_\discretionary{}{}{}MASKS\_\discretionary{}{}{}REST},
    mapped by (1+z\_s)/(1+z\_l)=1.374): Mg II 2796/2803 (def-rest
    3835--3855 Å), {[}O II{]} 3727/3729 (5115--5135), {[}Ne III{]} 3869
    (5260--5340), \textbf{He I 3819} (5237--5253, added 2026-05-27 as
    test M8).
  \item
    \textbf{Tellurics --- handled by masking, not correction.} One band
    is masked: the \textbf{O₂ A-band leading-edge residual} at
    \textbf{obs 7593--7626 Å = def-rest 4525--4545 Å}
    (\BreakablePath{ARC\_\discretionary{}{}{}MASKS\_\discretionary{}{}{}REST}).
    Confirmed telluric, not z=1.302 source, on 2026-05-18 by the
    same-observed-λ test (the spike sits at the \emph{same obs λ} in the
    deflector aperture, the sky box, and a 4--8″ off-deflector
    spectrum). No full telluric correction is applied; the M10 OH/sky
    bands (below) catch sky-emission residuals.
  \item
    \textbf{M10 sky-line audit:} bad-pixel catalog
    \BreakablePath{BAD\_\discretionary{}{}{}PIXELS\_\discretionary{}{}{}REST}
    = 35 entries (26 CR residuals + 9 OH airglow/sky bands), all
    verified non-source via the arc spectrum.
  \end{itemize}
\item
  \textbf{Wild bootstrap errors (→ paper §3.2.4; full detail in §2.2.3.4
  below).} Hybrid \textbf{Rademacher sign-flip × local-residual
  scaling}, N = 500 production per (SPS × degree), errors = 16/84
  percentiles. See §2.2.3.4 for the exact algorithm and the no-σ-clip
  note.
\end{itemize}

\hypertarget{bootstrap-error-analysis-paper-3.2.4}{%
\subsubsection{§2.2.3.4 Bootstrap error analysis (paper
§3.2.4)}\label{bootstrap-error-analysis-paper-3.2.4}}

\emph{This is the σ\_e \textbf{statistical} error. It is computed by a
hybrid wild bootstrap on the I-weighted R \textless{} R\_e aperture
spectrum, run independently for every (SPS library × additive polynomial
degree) and then pooled.
{[}\BreakablePath{scripts/\discretionary{}{}{}bootstrap\_\discretionary{}{}{}ppxf.\discretionary{}{}{}py}:
\BreakablePath{compute\_\discretionary{}{}{}local\_\discretionary{}{}{}residual\_\discretionary{}{}{}scaling},
\BreakablePath{run\_\discretionary{}{}{}bootstrap\_\discretionary{}{}{}single\_\discretionary{}{}{}degree},
\BreakablePath{run\_\discretionary{}{}{}bootstrap}{]}}

\begin{itemize}
\tightlist
\item
  \textbf{The algorithm, per (SPS × degree):}

  \begin{enumerate}
  \def\labelenumi{\arabic{enumi}.}
  \tightlist
  \item
    \textbf{Residuals} r = galaxy − best-fit, where best-fit is the
    original (un-resampled) ppxf fit at that degree.
  \item
    \textbf{Heteroskedastic local rescale (the ``local-residual''
    half).} A 75-pixel rolling window
    (\BreakablePath{uniform\_\discretionary{}{}{}filter1d},
    \BreakablePath{mode='reflect'}) gives the \textbf{local} residual
    std at each pixel; the scale factor = local\_std / global\_std,
    \textbf{clipped to {[}0.2, 5.0{]}}. scaled\_r = r × scale. This
    preserves the \textbf{wavelength-dependent noise structure} (sharp
    line cores vs smooth continuum) that a flat resample would erase.
  \item
    \textbf{Rademacher draw (the ``wild'' half).} Each pixel gets an
    independent sign s ∈ \{−1, +1\} with equal probability; the
    synthetic spectrum is \textbf{galaxy\_boot = best-fit + scaled\_r ×
    s}. The sign-flip keeps the pixel grid, the LOSVD line shape, and
    the local noise amplitude intact --- you cannot resample pixels
    independently here without destroying the line profile that σ is
    measured from.
  \item
    \textbf{Re-fit} ppxf on galaxy\_boot with the \emph{same}
    \BreakablePath{goodpixels}, \BreakablePath{moments=2},
    \BreakablePath{degree}, \BreakablePath{mdegree=0}; record V, σ, χ²,
    and z = (1+z₀)·exp(V/c) − 1.
  \end{enumerate}
\item
  \textbf{N = 500 production} per (SPS × degree); \textbf{N = 50 smoke}
  agrees to within 1 km/s. RNG \textbf{seeded}
  (\BreakablePath{np.\discretionary{}{}{}random.\discretionary{}{}{}default\_\discretionary{}{}{}rng(seed)})
  for reproducibility. Errors = \textbf{16th/84th percentiles}
  (asymmetric).
\item
  \textbf{Pooling = the systematic-marginalisation step.} All three SPS
  (FSPS+EMILES+XSL) × 15 polynomial degrees (15--29) bootstrap samples
  are \textbf{pooled into one posterior}. The pooled width therefore
  \textbf{already marginalises over the SPS-library and
  polynomial-degree systematics} --- between-SPS spread ±2.04,
  within-SPS bootstrap ±4.22, quadrature 4.73 ≈ pooled \textbf{4.64} →
  \textbf{no separate SPS or degree line is added to the budget (§3.1).}
\item
  \textbf{Headline stat error: ±4.6 km/s} (asymmetric \textbf{−5.16 /
  +4.13}) at the R\_e = 2.097″ aperture.
\item
  \textbf{No in-fit σ-clipping --- \BreakablePath{clean=False}} (the
  ppxf default; \BreakablePath{clean=} is \textbf{not} passed,
  \BreakablePath{bootstrap\_\discretionary{}{}{}ppxf.\discretionary{}{}{}py:434}).
  Rationale: \BreakablePath{clean=True} rejects \textbf{0 pixels} here
  because the KCWI noise array is \textbf{over-estimated} relative to
  the actual residual scatter, so iterative σ-clipping is a no-op.
  Outliers are instead removed \textbf{up front} by the explicit
  \textbf{35-entry
  \BreakablePath{BAD\_\discretionary{}{}{}PIXELS\_\discretionary{}{}{}REST}}
  mask (26 CR residuals + 9 OH/sky bands, all verified non-source via
  the arc spectrum), so every bootstrap iteration fits the \textbf{same
  fixed \BreakablePath{goodpixels} set} --- no per-iteration rejection
  that could bias σ.
\item
  \textbf{How the bad pixels were flagged (the reportable σ-clip
  number).}
  \BreakablePath{BAD\_\discretionary{}{}{}PIXELS\_\discretionary{}{}{}REST}
  was built by a \textbf{rolling 75-pixel local-MAD on the canonical
  ppxf-fit residuals, clipping at \textbar residual\textbar/local-MAD
  \textgreater{} 3σ} (rest 3800--5400 Å) → 52 outlier pixels in 26
  contiguous ranges (\textasciitilde46 Å, padded ±0.5 Å to survive
  log-rebinning); biggest outlier = a 6-pixel cluster at def-rest
  5232--5236 Å (obs 8768--8774 Å) at \textbf{26σ} in local-MAD units.
  \textbf{Intrinsic to the data} (same pixels in both
  \BreakablePath{\_\discretionary{}{}{}mtwdo\_\discretionary{}{}{}} and
  original reductions). The \textbf{M10 sky-line audit} then added
  \textbf{9 OH/airglow bands} by a separate threshold ---
  \textbf{\textgreater{} 2.5× the median sky std} --- all arc-verified
  non-source.
  {[}\BreakablePath{scripts/\discretionary{}{}{}run\_\discretionary{}{}{}window\_\discretionary{}{}{}sweep.\discretionary{}{}{}py:112–155}{]}
  \emph{Report as:} ``cosmic-ray residuals flagged at \textbf{3σ} in a
  75-pixel rolling local-MAD; sky/OH bands added above \textbf{2.5× the
  median sky noise} (35-entry catalog) --- distinct from the in-fit
  \BreakablePath{clean=False}.''
\item
  \textbf{Performance:} \BreakablePath{joblib} parallel, \textbf{BLAS
  pinned to 1 thread/worker} (avoids oversubscription);
  \textasciitilde12 min per SPS library at N=500. \textbf{Outputs:}
  \BreakablePath{results/\discretionary{}{}{}ppxf\_\discretionary{}{}{}bootstrap\_\discretionary{}{}{}errors\_\discretionary{}{}{}\{sps\}[\_\discretionary{}{}{}suffix].\discretionary{}{}{}npz},
  arrays shape (n\_degree, 500) for V/σ/χ²/z plus p16/p50/p84 summaries
  and asymmetric errors.
\item
  \textbf{Citation TODO (honest):} the repo cites \textbf{no}
  wild-bootstrap reference. The standard statistics precedents for the
  Rademacher-weight wild bootstrap are \textbf{Wu (1986), Mammen (1993),
  Davidson \& Flachaire (2008)} --- \emph{not yet in Zotero; verify +
  add before drafting, do not assert a specific one as the method source
  until checked.}
\end{itemize}

\textbf{σ\_e definition \& wavelength window}

\begin{itemize}
\item
  \textbf{σ\_e definition --- the luminosity-weighted second moment of
  the LOSVD within R\_e.} (Defined as the radial integral in
  \textbf{Gültekin et al.~2009, eq. (1)} {[}verified 2026-06-12{]}; the
  relation we place it on is \textbf{Kormendy \& Ho 2013 eq. 3}.) Two
  mathematically-equivalent representations:

  \begin{itemize}
  \item
    \textbf{Analytic / luminosity-weighted integral --- Gültekin et
    al.~2009 eq. (1)} (their 1D form; generalised to the 2D axisymmetric
    area element dA = 2πr dr):

    \[\sigma_e^2 = \frac{\int_0^{R_e} (V^2 + \sigma^2)\, I(r)\, 2\pi r\, dr}{\int_0^{R_e} I(r)\, 2\pi r\, dr}.\]
  \item
    \textbf{Discrete spaxel/bin sum} --- the same moment evaluated over
    the \(N\) spaxels (or Voronoi bins) \(i\) inside \(R_e\), with
    \(F_i\) = the I-band flux (luminosity weight) of spaxel \(i\); the
    2πr factor is captured implicitly because the number of spaxels per
    annulus scales with area:

    \[\sigma_e^2 = \frac{\sum_{i=1}^{N} F_i\,[(V_i - V_\mathrm{sys})^2 + \sigma_i^2]}{\sum_{i=1}^{N} F_i}.\]

    This discrete spaxel/bin sum \textbf{is Cappellari et al.~2013
    (ATLAS3D XV, MNRAS 432, 1709) eq. 29} ---
    \(\langle v_{\rm rms}^2\rangle_e \equiv \langle V^2+\sigma^2\rangle_e \equiv \sum_i F_i(V_i^2+\sigma_i^2)/\sum_i F_i\)
    --- equivalently the discretised form of the Gültekin 2009 eq. 1
    integral. (Their eq. 29 also carries an
    \textbf{inclination/projection element} in the dynamical-model
    framework; our §7 uses the \emph{observed-plane} projected
    \(V_n,\sigma_n\) directly and omits that term.) \textbf{It is NOT
    ``Cappellari 2006 eq. 1''} --- a prior draft mis-cited it there;
    Cappellari 2006 eq. 1 is the aperture correction
    \((\sigma_R/\sigma_e)=(R/R_e)^{-0.066}\), and its σ\_e (SAURON IV
    §2.3) is the co-add-then-fit method below. {[}eq. numbers verified
    2026-06-12: Cappellari 2006 eq. 1, Cappellari 2013 eq. 29, Gültekin
    2009 eq. 1{]}
  \end{itemize}
\item
  \textbf{Both forms are implemented and cross-checked here} (do not
  re-derive). The \textbf{headline} (§6cum/nb09) is the
  \textbf{co-add-then-fit method of Cappellari et al.~2006 (SAURON IV,
  §2.3)} --- co-add the I-weighted R\textless R\_e spectra into one
  effective spectrum and fit a single LOSVD with pPXF; the width of that
  co-added spectrum \emph{equals} the luminosity-weighted second moment
  above (this is what KH13/Greene+20/ATLAS\(^{\rm 3D}\) report). The
  \textbf{discrete annular sum (Gültekin 2009 eq. 1)} is the §7 / nb07c
  cross-check
  (\BreakablePath{run\_\discretionary{}{}{}gultekin\_\discretionary{}{}{}mc},
  F\_j = Σ I over each unmasked annulus). The two agree at
  \textbf{\textless1σ}: co-add-fit (§6cum, narrow) 267.32 ± 24 vs
  discrete annular (§7, arc-filtered to R\textless R\_safe=3R\_e/4)
  256.17 −13.0/+12.7. Full implementation + subtleties in
  \BreakablePath{reference\_\discretionary{}{}{}gultekin\_\discretionary{}{}{}implementation.\discretionary{}{}{}md}.
\item
  \textbf{★ What the headline code ACTUALLY computes}
  (\BreakablePath{final\_\discretionary{}{}{}sigma\_\discretionary{}{}{}e.\discretionary{}{}{}extract\_\discretionary{}{}{}aperture\_\discretionary{}{}{}spectrum}
  → pPXF; this is the number we report, 267.31). \textbf{We do NOT
  evaluate either σ\_e² sum or integral above.} We form the
  \textbf{I(r)-weighted co-added aperture spectrum}

  \[S_e(\lambda) \;=\; \sum_{i=1}^{N} w_i\, S_i(\lambda), \qquad
    w_i = \frac{I_i}{I_{\rm tot}}, \quad I_{\rm tot} \equiv \sum_{i=1}^{N} I_i, \qquad
    I_i \equiv 0 \ \text{for arc-masked spaxels (mask\_weight=0)},\]

  where \(i\) indexes the \(N\) spaxels inside \(R_e\) (\(r_i < R_e\)),
  \(S_i(\lambda)\) is the spaxel spectrum, and \(I_i\) its 6500--7500 Å
  observed-band IFU flux (the luminosity weight; arc-masked spaxels
  carry \(I_i=0\) so \(w_i=0\)). The denominator \(I_{\rm tot}\) is the
  single normalisation constant (total in-aperture flux) so
  \(\sum_i w_i = 1\). \textbf{The choice of weight map \(I_i\) is itself
  a budgeted systematic --- the ``I-shape'' term (±2.29 km/s; §2.4.3,
  \BreakablePath{run\_\discretionary{}{}{}isource\_\discretionary{}{}{}shape\_\discretionary{}{}{}sweep.\discretionary{}{}{}py}
  +
  \BreakablePath{ishape\_\discretionary{}{}{}arcfree\_\discretionary{}{}{}pooled.\discretionary{}{}{}py}).}
  Holding the spaxel selection fixed, we re-measure σ\_e with \textbf{14
  alternative weight maps}: the 10 of
  \BreakablePath{run\_\discretionary{}{}{}isource\_\discretionary{}{}{}shape\_\discretionary{}{}{}sweep}
  --- (1) the headline 6500--7500 Å IFU band,

  \begin{enumerate}
  \def\labelenumi{(\arabic{enumi})}
  \setcounter{enumi}{1}
  \tightlist
  \item
    full IFU white-light, (3--4) HST F140W / F200LP reprojected (raw),
    (5--6) the same arc-masked, (7--8) their 1-D curve-of-growth annular
    means, (9--10) their 2-D Sérsic-model fits --- \textbf{plus 4
    arc-free PSF-matched maps} (11--14: F140W/F200LP single-Sérsic and
    arc-filled, each convolved to the 1.27″ KCWI PSF; see below).
    Peak-to-peak/2 of the full 14-shape spread (266.79--271.37 km/s, NEW
    cube + M10) = \textbf{±2.29 km/s}. So \(I_i\) = IFU white-light is
    the \emph{headline} choice, not an assumption: σ\_e moves ≤±2.29
    km/s across this whole family of luminosity weightings.
  \end{enumerate}

  \begin{itemize}
  \tightlist
  \item
    \textbf{PSF resolution of the HST weight maps (FOLDED into I-shape
    2026-06-12).} The HST F140W/F200LP (and Sérsic) weight maps are
    reprojected at native HST resolution, not convolved to the KCWI
    \textasciitilde1.27″ seeing. To test resolution-matching
    \textbf{correctly the arc must be removed first} --- using the
    \emph{arc-free deflector model} (single-Sérsic, or the HST image
    with the masked arc filled by that model) and only \emph{then}
    convolving to the KCWI PSF
    (\BreakablePath{scripts/\discretionary{}{}{}ishape\_\discretionary{}{}{}arcfree\_\discretionary{}{}{}psf\_\discretionary{}{}{}test.\discretionary{}{}{}py},
    \BreakablePath{ishape\_\discretionary{}{}{}arcfree\_\discretionary{}{}{}pooled.\discretionary{}{}{}py}).
    On the 3-SPS-pooled footing these arc-free, PSF-matched weights give
    σ\_e = \textbf{266.79--268.78} km/s (F140W Sérsic 266.79 / filled
    267.38; F200LP Sérsic 268.29 / filled 268.78) --- i.e.~they fall
    \textbf{inside} the existing 10-shape spread (266.83--271.37).
    Folding them in (14 shapes) moves the I-shape term \textbf{±2.27 →
    ±2.29} (negligible) and σ\_e is unchanged. \emph{(Cautionary
    cross-check: naively convolving the \textbf{raw} arc-containing HST
    images instead gives a spurious +7--8 km/s --- that is arc
    \textbf{leakage} (the 1.27″ kernel spreads the masked arc light into
    unmasked arc-adjacent spaxels, worst for the UV-bright F200LP) plus
    high-σ-outskirt up-weighting, NOT a resolution effect; not used.)}
    The \textbf{headline IFU white-light} weight is already at KCWI
    resolution AND arc-masked AND deflector-light- weighted, so the
    resolution choice is bounded within the I-shape term. pPXF then fits
    a \textbf{single Gaussian LOSVD} (\BreakablePath{moments=2}) to
    \(S_e(\lambda)\),
  \end{itemize}

  \[S_e(\lambda) \;\approx\; \Big[\textstyle\sum_k a_k\, T_k(\lambda)\Big] \otimes \mathcal{G}(v;\,V,\sigma)
    \;+\; \text{(additive Legendre poly, deg 15–29)},\]

  and \textbf{\(\sigma_e \equiv \sigma\)}, the fitted LOSVD dispersion.
  Because superposing spaxel spectra at different line-of-sight
  velocities \(V_i\) broadens the co-add, the width \(\sigma\) of that
  single fit \emph{equals} the luminosity-weighted second moment
  \(\sqrt{\langle (V_i-V_{\rm sys})^2 + \sigma_i^2\rangle}\) --- so it
  \textbf{measures} the Cappellari 2013 eq. 29 / Gültekin 2009 eq. 1
  quantity \textbf{without ever forming the explicit sum} (and with
  \(V_{\rm sys}\) absorbed automatically into the fitted \(V\); §2.4.1).
  Code:
  \BreakablePath{extract\_\discretionary{}{}{}aperture\_\discretionary{}{}{}spectrum}
  (the weighted sum) +
  \BreakablePath{bootstrap\_\discretionary{}{}{}ppxf.\discretionary{}{}{}setup\_\discretionary{}{}{}ppxf\_\discretionary{}{}{}inputs\_\discretionary{}{}{}from\_\discretionary{}{}{}spectrum}
  +
  \BreakablePath{ppxf(.\discretionary{}{}{}.\discretionary{}{}{}.\discretionary{}{}{}, moments=2)}.
  The explicit discrete sum \(\sum_j F_j(V_j^2+\sigma_j^2)/\sum_j F_j\)
  is computed \textbf{only} in the §7 annular cross-check
  (\BreakablePath{run\_\discretionary{}{}{}gultekin\_\discretionary{}{}{}mc}),
  never in the headline.
\item
  \textbf{Why the integral path is the headline} (not the discrete sum):
  the discrete spaxel-sum is (i) per-spaxel-S/N-limited and (ii)
  sensitive to arc contamination in the outer spaxels --- an
  \emph{unfiltered} discrete sum over the resolved PowerBins out to
  \textasciitilde3.4″ is dominated by the z=1.302 arc (outer bins carry
  σ ≳ 300--440 km/s, V swings of ±150--270; the σ\_e blows up to
  \textasciitilde350). The I-weighted aperture auto-suppresses the arc
  (bright deflector core dominates the weight) and avoids the S/N floor,
  so it implements the same Cappellari definition robustly. The discrete
  §7 sum is therefore arc-filtered to R\textless R\_safe and used only
  as the architecture-independent cross-check / for the σ(R) profile.
\item
  \textbf{Standard wavelength range (headline): rest 3800--5400 Å}
  (\BreakablePath{wR3800\_\discretionary{}{}{}5400\_\discretionary{}{}{}arcmask})
  --- Ca H+K through Mg b / Fe5270 / Fe5335; 2161 good pixels
  (\textasciitilde2.6× the narrow window). The fit is constrained by the
  highest-S/N features --- \textbf{Ca H, Ca K, G-band} (3800--4400) ---
  and is \textbf{consistent across windows}.

  \begin{itemize}
  \tightlist
  \item
    Blue edge 3800 Å = XSL template floor (def-rest 3675 Å) + 5\%
    velocity padding + retaining Ca H+K; red edge set by KCWI coverage.
    \textbf{Provenance is data-driven, NOT a specific paper.} Closest
    precedent: LEGA-C (Bezanson et al.~2018; van der Wel et al.~2021) at
    z\textasciitilde0.7. \textbf{DO NOT cite ``Cappellari 2017
    recommends extending blueward'' --- that was a hallucination,
    retracted.}
  \item
    Cross-check window: \textbf{rest 4000--5400 Å}
    (\BreakablePath{wR4000\_\discretionary{}{}{}5400\_\discretionary{}{}{}arcmask},
    Hβ + Mg b, \textbf{no Ca H+K}) --- orthogonal Lick-tradition feature
    set.
  \item
    Narrow legacy window: obs 6500--7500 Å = rest 3879--4476 Å (Ca H+K +
    G); σ\_e = 267.95 ± 30.10.
  \end{itemize}
\end{itemize}

\textbf{A few systematics tests to describe} (full budget in Results
§3.1): I-shape (10 I-weight-map shapes, ±2.27), F200 mask weight
(w∈\{0,0.5,1\}, ±6.65), fit-window (3 windows, ±3.82), frame (vac/air,
±5.0), centering (5-center ±0.4″ sweep, ±4.0), reduction-pass (NEW vs
OLD cube, ±3.45). Plus the \textbf{M9 ``DO NOT MASK'' finding}: the
+5--7\% bump at def-rest 5193--5204 Å is the Mg b LOSVD wing (signal,
not contamination); masking it drops σ\_e by 7 km/s.

\begin{center}\rule{0.5\linewidth}{0.5pt}\end{center}

\hypertarget{lens-modeling}{%
\subsection{§2.3 Lens Modeling}\label{lens-modeling}}

\begin{quote}
\textbf{G1 --- now sourced from the companion-paper draft (2026-06-12).}
The lens-modeling content below is summarised from the \textbf{Ferrami
et al.~draft} (G. Ferrami, MPA Garching; lead-author email
gferrami@mpa-garching.mpg.de), \emph{``A supermassive black hole in the
Einstein spiral DES J0206−0114''}, file
\BreakablePath{\~/\discretionary{}{}{}Documents/\discretionary{}{}{}AGEL/\discretionary{}{}{}AGEL\_\discretionary{}{}{}0206\_\discretionary{}{}{}ApJL\_\discretionary{}{}{}Figures/\discretionary{}{}{}SMBH\_\discretionary{}{}{}in\_\discretionary{}{}{}Einstein\_\discretionary{}{}{}Spiral\_\discretionary{}{}{}DESJ0206\_\discretionary{}{}{}0114\_\discretionary{}{}{}June12/\discretionary{}{}{}main.\discretionary{}{}{}md}.
\textbf{This is a live draft with explicit
\BreakablePath{PLACEHOLDER UNTIL FINAL RUNS ARE COMPLETED} flags on the
evidence table, the M•--σ figure, and the potential-correction figure
--- every numeric BH/evidence value here is provisional.} The
complementary on-hand pointer: PROJECT\_BRIEF.md notes the lens-model
enclosed mass combines with (σ, M⋆, z) to give a 4D constraint on the
central BH + bulge scaling relations. \emph{This is OUR
(kinematics+photometry) paper; the lens model is the companion paper ---
cite it, do not reproduce its results as ours.}
\end{quote}

\textbf{⚠ Cross-paper reconciliation flags (resolve before any joint
number is quoted):}

\begin{itemize}
\tightlist
\item
  \textbf{Cosmology --- RESOLVED 2026-06-12: we adopted Planck 2015 to
  match Ferrami} (Planck Collaboration 2016): \textbf{H₀ = 67.7
  km/s/Mpc, Ω\_m,0 = 0.302}. (Previously our paper used H₀=70,
  Ω\_m=0.3.) Impact on our numbers is small and documented in the
  headline ``Cosmology switch --- impact'' box: R\_e 14.76→15.26 kpc,
  log M⋆ 11.47→11.50 (+0.028 dex, exact D\_L² rescale), σ\_e and all
  angular quantities unchanged. Now both papers share a cosmology → a
  Ferrami M\_• and our M⋆/σ\_e can be co-plotted directly. The Ferrami
  draft still carries its own standing TODO \emph{``Check all mass
  values to be corrected for the different deflector redshift''} ---
  that is the redshift (not cosmology) reconciliation below, and our z
  is the canonical one.
\item
  \textbf{Redshifts differ slightly.} Ferrami quotes \textbf{z\_l =
  0.675, z\_s = 1.303}; our systemic values (notebook 04) are
  \textbf{z\_l = 0.67564, z\_s = 1.30263}. Minor, but the σ\_inst/V\_sys
  frame and the M• conversion both use z\_l --- reconcile to our
  line-fit values when combining.
\item
  \textbf{System identity (agreed):} \textbf{AGEL J020613−011417} = DES
  J0206−0114, AGEL DR2 (Barone et al. 2026); ICRS RA = 31.55611°, Dec =
  −1.23817°. Discovered in DES via a CNN (Jacobs et al.). HST WFC3
  \textbf{F200LP + F140W}, SNAP program \textbf{16773 (PI Glazebrook)},
  two 300 s exposures per filter --- the same imaging this paper uses
  for photometry. The lens forms a \textbf{radial arc} (three
  star-forming source galaxies near a fold + a cusp); the shallow,
  sub-isothermal total slope is what makes the radial image bright.
\end{itemize}

\hypertarget{modelling-code-pronto}{%
\subsubsection{\texorpdfstring{§2.3.1 Modelling code ---
\BreakablePath{pronto}}{§2.3.1 Modelling code --- }}\label{modelling-code-pronto}}

\begin{itemize}
\tightlist
\item
  \textbf{\BreakablePath{pronto}} (Vegetti \& Koopmans 2009; Rybak et
  al.~2015; Rizzo et al.~2018; Ritondale et al.~2019; Powell et
  al.~2020; Ndiritu et al.~2025). Simultaneously recovers a
  \textbf{resolved background source} (regularised on a \textbf{Delaunay
  adaptive tessellation}) and the non-linear parameters of the smooth
  deflector \textbf{light + mass} distributions. The smooth potential
  can additionally be perturbed by \textbf{regularised pixellated
  potential corrections} (a.k.a. \emph{gravitational imaging}).
  Hierarchical Bayesian framework; \textbf{evidence ln 𝓔} computed by
  exploring the posterior with \textbf{MultiNest} (Feroz et al.;
  importance-nested-sampling variant). {[}Ferrami draft §Method{]}
\end{itemize}

\hypertarget{mass-density-profiles}{%
\subsubsection{§2.3.2 Mass density
profiles}\label{mass-density-profiles}}

\begin{itemize}
\tightlist
\item
  \textbf{Main lens = main elliptical parametric profile + mass
  multipoles + 2 secondary SIS deflectors + external shear.} Three trial
  forms for the main elliptical component:

  \begin{itemize}
  \tightlist
  \item
    \textbf{sEPL} --- singular elliptical power-law (the cEPL limit r\_c
    → 0). \emph{This is the main halo of the best model.}
  \item
    \textbf{cEPL} --- cored elliptical power-law: κ(ξ) =
    κ₀·(4−γ)/(4√q)·(r\_c² + ξ²)\^{}((1−γ)/2), ξ = √(x²+y²/q²) the
    elliptical radius, r\_c the core radius, γ the 3D slope, q the axis
    ratio. Deflection via \textbf{\BreakablePath{fastell}} (Barkana).
    Einstein-radius definition for the cEPL from \textbf{Enzi et
    al.~(cEPL θ\_E)}.
  \item
    \textbf{bEPL} --- broken elliptical power-law: piecewise κ with
    break radius θ\_B, break convergence κ\_B, inner/outer 2D slopes
    t₁/t₂ (O'Riordan, BPL derivation).
  \end{itemize}
\item
  \textbf{Two companion galaxies} within \textasciitilde4″ of the main
  lens → \textbf{singular isothermal spheres (SIS)} (cEPL with q=1, γ=2,
  r\_c=0).
\item
  \textbf{Multipoles of order m = 1, 3, 4}, \emph{circular
  slope-matched}: κ(θ,φ) = θ\^{}(1−γ)·{[}a\_m sin(mφ) + b\_m cos(mφ){]}.
  Choice of circular (not elliptical-isothermal) multipoles justified
  via Paugnat et al.

  \begin{enumerate}
  \def\labelenumi{(\arabic{enumi})}
  \setcounter{enumi}{2024}
  \tightlist
  \item
    eq. B10: at the recovered γ≈1.3, q≈0.7, circular multipoles are the
    better approximation for θ outside the band (1 ± 0.15)·θ\_E --- and
    the science target is the BH perturbation at \textbf{θ ≪ θ\_E}.
  \end{enumerate}
\item
  \textbf{External shear} of strength Γ and PA θ\_Γ.
\end{itemize}

\hypertarget{central-perturber-the-bh-test}{%
\subsubsection{§2.3.3 Central perturber (the BH
test)}\label{central-perturber-the-bh-test}}

\begin{itemize}
\tightlist
\item
  After fixing the smooth-model parameters, the draft tests whether the
  data \textbf{require an extra concentrated central mass} by comparing
  Bayesian evidence across perturber choices:

  \begin{itemize}
  \tightlist
  \item
    \textbf{Point mass (PM)} --- κ(θ) =
    (θ\_\{E,•\}²/2)·θ/\textbar θ\textbar², θ\_\{E,•\} the Einstein
    radius of a point mass M\_• at z\_l. Run \textbf{twice}: (i)
    \textbf{fixed} to the centre of the foreground mass, and (ii)
    \textbf{free} within a \textbf{0.5″} radius of the light centre.
  \item
    \textbf{SPL} --- spherical power-law (EPL with q=1, r\_c=0).
  \item
    \textbf{NFW} --- Navarro--Frenk--White with \textbf{free
    concentration c} (κ\_s = ρ\_s r\_s Σ\_c). The PM is the SMBH
    hypothesis; SPL/NFW are the extended-mass (dark-subhalo)
    alternatives it is tested against. {[}Ferrami draft §Method →
    Central black hole{]}
  \end{itemize}
\end{itemize}

\hypertarget{lens-light-source-regularisation}{%
\subsubsection{§2.3.4 Lens light + source
regularisation}\label{lens-light-source-regularisation}}

\begin{itemize}
\tightlist
\item
  \textbf{Lens light = a single Sérsic with m = 1,3,4 multipole
  perturbations}, fit \emph{simultaneously} with the lens + BH mass in
  \BreakablePath{pronto}. Evidence comparison disfavoured both a 2nd/3rd
  Sérsic and the no-multipole model. During burn-in, \textbf{3 sets of
  multiple images} are required to focus to a tolerance shrinking to
  \textbf{0.05″}.
\item
  \textbf{Source = Delaunay adaptive tessellation with adaptive,
  SNR-weighted regularisation} (O'Riordan, Euclid lens): two forms
  tested --- penalising the \textbf{gradient} vs the \textbf{curvature}
  of the source surface brightness; per-pixel weight log(λ\_i/λ\_s) =
  −w\_a·x\_i\^{}(−w\_b), x\_i the log-SNR normalised over all pixels,
  with w\_a, w\_b free. \textbf{Curvature regularisation wins by Δln 𝓔 ≈
  100} over gradient → the paper uses \textbf{curvature} throughout.
  {[}Ferrami draft §Adaptive source regularisation{]}
\end{itemize}

System identity for the summary sentence: \textbf{AGEL J020613−011417}
(DES J0206−01), AGEL DR2; source z = 1.302 (spiral), deflector z =
0.67564.

\begin{center}\rule{0.5\linewidth}{0.5pt}\end{center}

\hypertarget{end-to-end-procedures-ordered-walkthroughs}{%
\subsection{§2.4 End-to-end procedures (ordered
walkthroughs)}\label{end-to-end-procedures-ordered-walkthroughs}}

\emph{This section spells out the actual step order of each pipeline ---
what runs, in what sequence, with which script --- so the Methods
section and a referee can reconstruct every headline number. Numeric
results are the M12 headline
(\BreakablePath{results/\discretionary{}{}{}PAPER\_\discretionary{}{}{}VALUES.\discretionary{}{}{}json}).}

\hypertarget{spectroscopic-ux3c3_e-pipeline-step-by-step}{%
\subsubsection{§2.4.1 Spectroscopic σ\_e pipeline --- step by
step}\label{spectroscopic-ux3c3_e-pipeline-step-by-step}}

Driver for the headline:
\BreakablePath{scripts/\discretionary{}{}{}run\_\discretionary{}{}{}wide\_\discretionary{}{}{}sigma\_\discretionary{}{}{}e.\discretionary{}{}{}py -\discretionary{}{}{}-\discretionary{}{}{}cube new\_\discretionary{}{}{}clean\_\discretionary{}{}{}hei -\discretionary{}{}{}-\discretionary{}{}{}n\_\discretionary{}{}{}bootstrap 500}
(calls
\BreakablePath{final\_\discretionary{}{}{}sigma\_\discretionary{}{}{}e.\discretionary{}{}{}load\_\discretionary{}{}{}setup}
→
\BreakablePath{extract\_\discretionary{}{}{}aperture\_\discretionary{}{}{}spectrum}
→ \BreakablePath{bootstrap\_\discretionary{}{}{}ppxf} machinery).

\begin{enumerate}
\def\labelenumi{\arabic{enumi}.}
\tightlist
\item
  \textbf{Reduction (upstream).} Keck II / KCWI, KCRM red arm, Red-Low
  grating, Medium slicer, 2×2 binning, dithered PA 0/45/90°. The
  headline cube is the NEW
  \textbf{\BreakablePath{\_\discretionary{}{}{}mtwdo\_\discretionary{}{}{}}}
  reduction (Master-Twilight+Dome hybrid flats, K. R. Gupta),
  \BreakablePath{Nov17\_\discretionary{}{}{}2025\_\discretionary{}{}{}DESJ0206\_\discretionary{}{}{}RL\_\discretionary{}{}{}combined\_\discretionary{}{}{}mtwdo\_\discretionary{}{}{}icubes\_\discretionary{}{}{}wcs.\discretionary{}{}{}fits}.
  The
  pre-\BreakablePath{\_\discretionary{}{}{}mtwdo\_\discretionary{}{}{}}
  cube is the OLD reduction (second-reduction cross-check). Cube axes
  \BreakablePath{[λ, y, x]}; 0.30″/spaxel; seeing FWHM 1.27″
  (\BreakablePath{GUIDFWHM}); σ\_inst ≈ 12.6 km/s; R ≈ 10000.
\item
  \textbf{Geometry + R\_e.} HST-mean centre =
  \BreakablePath{photutils.\discretionary{}{}{}centroid\_\discretionary{}{}{}2dg}
  of the F140W+F200LP cores (offset 0.36″; F200LP pulled by the UV arc,
  so F140W is the clean bulge centroid), propagated through the KCWI WCS
  to a sub-pixel IFU centre. R\_e = mean of the F140W (1.912″) + F200LP
  (2.281″) \textbf{best-mask} curves-of-growth = \textbf{2.097″}
  (\BreakablePath{final\_\discretionary{}{}{}sigma\_\discretionary{}{}{}e.\discretionary{}{}{}curve\_\discretionary{}{}{}of\_\discretionary{}{}{}growth},
  r\_max=6″; photutils-validated ±0.002″).
\item
  \textbf{Spatial arc mask → IFU grid.} The F200LP
  \BreakablePath{\_\discretionary{}{}{}mask.\discretionary{}{}{}fits}
  (2512 HST px, arc-only) is reprojected onto the 0.30″ IFU grid with
  \BreakablePath{scipy.\discretionary{}{}{}ndimage.\discretionary{}{}{}map\_\discretionary{}{}{}coordinates(order=0)}
  (nearest-neighbour) →
  \BreakablePath{arc\_\discretionary{}{}{}spax\_\discretionary{}{}{}mask}.
  \textbf{Mask coverage within R\_e (PSF caveat;
  \BreakablePath{scripts/\discretionary{}{}{}psf\_\discretionary{}{}{}mask\_\discretionary{}{}{}fraction.\discretionary{}{}{}py}):}
  this \emph{geometric} mask flags \textbf{35/154 = 23\% of spaxels
  (28\% by I-weight)} inside R\_e=2.097″. But the arc light is spread by
  the \textbf{1.27″ KCWI seeing}, so the geometric fraction is a
  \textbf{lower bound}: growing the mask by the PSF (binary dilation)
  raises the masked fraction to \textbf{71\% / 81\%} (by FWHM/2 = 0.64″)
  or \textbf{96\% / 98\%} (by the full FWHM = 1.27″) --- i.e.~a hard
  seeing-grown mask removes nearly the entire aperture and is
  \textbf{infeasible}. (Convolving the binary mask with the PSF and
  cutting at \textgreater25\% / \textgreater10\% arc-light contamination
  gives 42\% / 75\% by count; the \textgreater50\% cut, 1\%, is a
  thin-arc dilution artifact.) This is precisely why we instead (i)
  \textbf{I-weight} the aperture --- the bright deflector core
  dominates, the arc-contaminated outskirts are down-weighted --- and
  (ii) carry the \textbf{mask-weight systematic (±6.65 km/s, w=0→1;
  §2.4.3)} to bracket the residual arc dilution rather than attempt a
  hard PSF-grown mask.
\item
  \textbf{I-weighted aperture extraction.} Build the 6500--7500 Å
  white-light I-map; the R\textless R\_e aperture spectrum = Σ\_spaxel
  cube · w, w = I-weight with arc spaxels \textbf{hard-masked
  (mask\_weight=0)} and re-normalised. Spectra are summed (not averaged)
  so the bootstrap preserves per-spaxel noise.
  (\BreakablePath{extract\_\discretionary{}{}{}aperture\_\discretionary{}{}{}spectrum},
  \BreakablePath{final\_\discretionary{}{}{}sigma\_\discretionary{}{}{}e.\discretionary{}{}{}py}.)
\item
  \textbf{Spectral masking (goodpixels).} (a) \textbf{no-Balmer} mask
  --- forbidden lines only ({[}O II{]}, {[}O III{]}, {[}O I{]}, {[}N
  II{]}, {[}S II{]}); \textbf{Hδ, Hγ, Hβ kept} (absorption; M12/nb16
  confirmed keep-unmasked).

  \begin{enumerate}
  \def\labelenumii{(\alph{enumii})}
  \setcounter{enumii}{1}
  \tightlist
  \item
    \textbf{z=1.302 source-emission} masks
    \BreakablePath{ARC\_\discretionary{}{}{}MASKS\_\discretionary{}{}{}REST}
    (Mg II, {[}O II{]}, {[}Ne III{]}, \textbf{He I 3819}) + the O₂
    A-band telluric. (c)
    \textbf{\BreakablePath{BAD\_\discretionary{}{}{}PIXELS\_\discretionary{}{}{}REST}}
    = 35 entries (26 CR residuals + 9 M10 OH/sky bands). All in
    deflector rest-frame Å.
  \end{enumerate}

  \begin{itemize}
  \tightlist
  \item
    \textbf{Outlier rejection / sigma-clipping decision:} ppxf's in-fit
    \textbf{\BreakablePath{clean=False}} --- \BreakablePath{clean=True}
    rejects 0 pixels here (the KCWI noise array is overestimated vs the
    actual residual scatter), so it is inert. Instead, cosmic-ray/bad
    pixels are flagged by a \textbf{3σ local-MAD clip} (rolling 75-pixel
    window, \textbar residual\textbar/local-MAD \textgreater{} 3σ) on
    the canonical residuals → 52 CR-like spikes, curated into the 26 CR
    entries of
    \BreakablePath{BAD\_\discretionary{}{}{}PIXELS\_\discretionary{}{}{}REST};
    the 9 M10 sky bands were flagged separately at \textgreater2.5×
    median sky-noise. \textbf{Caveats (carried as a stated sensitivity,
    not a budget term):} the 3σ threshold is not strongly motivated vs
    alternatives, and cleaning shifts σ\_e by only +1.4 km/s (below stat
    1σ); the same pixels replicate the shift on the OLD cube (M6) → real
    CR residuals.
  \end{itemize}
\item
  \textbf{Per-SPS frame-aware ppxf.} For each of FSPS / EMILES / XSL:
  match the galaxy frame to the SPS native frame (FSPS=vacuum,
  EMILES/XSL=air; scalar-median vac↔air via Ciddor 1996), run
  \BreakablePath{ppxf} (Cappellari) with \textbf{moments=2},
  \textbf{additive Legendre \BreakablePath{degree} swept 15--29 (15
  values)}, \BreakablePath{mdegree=0}, over the wide window (rest
  3800--5400 Å). Subtract the per-SPS V\_sys before pooling.
\item
  \textbf{Wild-bootstrap errors.} Per (SPS × degree), \textbf{N=500}
  hybrid wild bootstrap
  (\BreakablePath{bootstrap\_\discretionary{}{}{}ppxf.\discretionary{}{}{}compute\_\discretionary{}{}{}local\_\discretionary{}{}{}residual\_\discretionary{}{}{}scaling}):
  residuals r = galaxy − bestfit; a rolling \textbf{75-pixel} window
  estimates the local (wavelength-dependent) noise scale s(λ) (clipped
  to \textbf{{[}0.2, 5.0{]}}); each iteration applies a
  \textbf{Rademacher} ±1 sign-flip to the locally-rescaled residual,
  adds it back to the best fit, and re-fits → one σ draw. Seeds are
  deterministic
  (\BreakablePath{BOOT\_\discretionary{}{}{}SEED=42 + 50000 + 10000·W\_\discretionary{}{}{}IDX + 100·sps\_\discretionary{}{}{}idx + degree\_\discretionary{}{}{}i})
  so caches bit-reproduce. joblib \BreakablePath{loky}, \textbf{BLAS
  pinned to 1 thread/worker}; N=50 smoke before N=500 production.
\item
  \textbf{Pool + budget.} Concatenate all σ samples (3 SPS × 15 degrees
  × 500) → 16/50/84 percentiles = central + asymmetric stat. The pooled
  width already marginalises SPS + polynomial degree. Add the §2.4.3
  systematic components in quadrature → \textbf{σ\_e = 267.31
  −12.98/+12.61 (sym ±12.79) km/s} (at the best-mask R\_e=2.097″
  aperture; cache
  \BreakablePath{results/\discretionary{}{}{}run\_\discretionary{}{}{}wide\_\discretionary{}{}{}sigma\_\discretionary{}{}{}e/\discretionary{}{}{}resys\_\discretionary{}{}{}best\_\discretionary{}{}{}mean/\discretionary{}{}{}}).
\end{enumerate}

\hypertarget{photometry-m-pipeline-step-by-step}{%
\subsubsection{§2.4.2 Photometry / M⋆ pipeline --- step by
step}\label{photometry-m-pipeline-step-by-step}}

Drivers:
\BreakablePath{scripts/\discretionary{}{}{}photometry\_\discretionary{}{}{}systematics.\discretionary{}{}{}py}
(masking + photometry),
\BreakablePath{bagpipes\_\discretionary{}{}{}sersic\_\discretionary{}{}{}refit.\discretionary{}{}{}py}
(SED), orchestrated/displayed in notebook 12. Recipe = \textbf{F200LP
locates, IR extends, fill-in recovers}.

\begin{enumerate}
\def\labelenumi{\arabic{enumi}.}
\tightlist
\item
  \textbf{Locate the arc on F200LP} (best source contrast): a 2D Sérsic
  deflector model is subtracted and a pixel is flagged as arc where its
  positive residual exceeds \textbf{k·σ\_sky} --- i.e.~\textbf{k is the
  detection threshold in units of the sky noise σ\_sky} (so k=3 means
  residual \textgreater{} 3σ\_sky), with a parallel S/N floor (the
  pixel's own S/N \textgreater{} SNR\_thresh). We adopt \textbf{k = 3}
  (\BreakablePath{arc\_\discretionary{}{}{}mask\_\discretionary{}{}{}verification.\discretionary{}{}{}py}):
  this \textbf{objective Sérsic-residual mask reproduces the expert hand
  mask to 0.016 mag}, R\_e to 3\%. \textbf{k = 3 is the clean saturation
  point} of a 2-D sweep over the two thresholds (SNR\_thresh ∈ \{2..20\}
  × k ∈ \{2..8\}) --- below k≈3 it grabs noise, above it the masked-flux
  fraction stops changing.
\item
  \textbf{Reproject the F200LP footprint to every band} (WCS +
  \BreakablePath{map\_\discretionary{}{}{}coordinates}) --- reproduces
  the expert JWST aperture mags to 0.01--0.02 mag.
\item
  \textbf{IR-extend with a COLOR GATE.} In the deeper/sharper JWST
  bands, region-grow the arc seed into pixels that are \textbf{(a)}
  significant positive residual above the single-Sérsic deflector model
  (data−model \textgreater{} 3σ\_sky), \textbf{(b)} bluer than the
  deflector core by ≥0.5 mag in F200LP−F140W, and \textbf{(c)}
  contiguous with the seed within R\textless5″. The \textbf{color gate
  (b) is essential}: a single Sérsic slightly under-fits the bright IR
  body, so residual-only growth would sweep the \emph{red} galaxy body
  and companions into the mask; the \emph{blue} z=1.302 source vs
  \emph{red} passive deflector contrast keeps only genuine source
  pixels. See §2.1.1b for the detailed rationale + caveats.
\item
  \textbf{Deflector model = single-component Sérsic} --- physically
  motivated for the elliptical (per user 2026-06-10); not a bulge+disk
  decomposition. Fit gotcha:
  \BreakablePath{Sersic2D.\discretionary{}{}{}amplitude} is I(r\_eff)
  not the peak --- init \textasciitilde peak·exp(−b\_n), bound to
  (peak·1e-3, peak·1.5), else LevMar collapses to ≈0. PSF-convolved fill
  quantified at ≤0.004 mag
  (\BreakablePath{psf\_\discretionary{}{}{}fill\_\discretionary{}{}{}model.\discretionary{}{}{}py})
  --- negligible.
\item
  \textbf{Two photometry estimators per band/mask:} \textbf{raw}
  \BreakablePath{photutils} aperture (masked → discards under-arc
  deflector light → biased low, mask-size-dependent) vs
  \textbf{single-Sérsic fill-in} (replace masked pixels with the model →
  recovers it). Fill-in correction \textbf{−0.19 to −0.27 mag} under the
  color-gated masks (smaller than the old 2-component +0.18--0.96 mag
  --- cleaner masks remove less).
\item
  \textbf{AB magnitudes} from per-image header zeropoints (never
  hardcoded): HST \BreakablePath{PHOTFLAM/\discretionary{}{}{}PHOTPLAM},
  JWST \BreakablePath{PIXAR\_\discretionary{}{}{}SR}. Four bands:
  F200LP, F140W (HST WFC3), F150W2, F322W2 (JWST NIRCam).
\item
  \textbf{Bagpipes SED}
  (\BreakablePath{02\_\discretionary{}{}{}streamlined\_\discretionary{}{}{}Bagpipes\_\discretionary{}{}{}SED}):
  exponential-τ SFH, Calzetti dust (A\_V 0--2), free Z (0--2.5 Z☉), z
  prior (0.674,0.676); \textbf{Nautilus} sampler n\_live=400, 500-sample
  posterior. A \textbf{10\% fractional flux floor} per band covers
  photon noise + the drizzle pixel-correlation under-estimate (empirical
  noise ≈1.3--2× CCD-equation); \textbf{20\%} is run as a conservative
  variant.
\item
  \textbf{M⋆ budget} = 8 fits \{per-band, global\} × \{raw, filled\} ×
  \{10\%, 20\%\} (script-reproducible via
  \BreakablePath{scripts/\discretionary{}{}{}mstar\_\discretionary{}{}{}masking\_\discretionary{}{}{}budget.\discretionary{}{}{}py}
  →
  \BreakablePath{Mstar\_\discretionary{}{}{}masking\_\discretionary{}{}{}systematic.\discretionary{}{}{}npz},
  \BreakablePath{photometry\_\discretionary{}{}{}systematics\_\discretionary{}{}{}Mstar.\discretionary{}{}{}npz})
  → headline (empirical, raw-central with the best per-band color-gated
  mask, one-sided +sys, user choice): \textbf{log M⋆ = 11.36 ± 0.08
  (stat) +0.14 (sys) at 10\%} {[}11.26 ± 0.14 +0.18 at 20\%{]}, fill-in
  reach 11.47; explicit masking-approach systematic \textbf{±0.086 dex}
  (was ±0.16 under the 2-component masks). Central consistent with the
  validated expert hand-mask 11.33.
\end{enumerate}

\hypertarget{systematic-audit-procedures-how-each-budget-term-was-measured}{%
\subsubsection{§2.4.3 Systematic-audit procedures (how each budget term
was
measured)}\label{systematic-audit-procedures-how-each-budget-term-was-measured}}

Each σ\_e systematic is a \textbf{dedicated sweep} re-using the §2.4.1
machinery, varying ONE axis and quoting \textbf{peak-to-peak/2} (or
half-Δ for two-point axes), all on the headline cube + masks:

\begin{itemize}
\tightlist
\item
  \textbf{I-shape (±2.27):} 10 alternative I-weight maps
  (\BreakablePath{run\_\discretionary{}{}{}isource\_\discretionary{}{}{}shape\_\discretionary{}{}{}sweep.\discretionary{}{}{}py})
  × 3 SPS × 15 deg × N=250; peak-to-peak/2 of the 10 central σ\_e
  (266.83--271.37).
\item
  \textbf{F200 mask-weight (±6.65):} down-weight the arc spaxels at
  w∈\{0,0.5,1\} × 3 SPS × N=500; (w00 269.69 − w100 256.39)/2.
  \emph{(This is the mask-\textbf{strength} axis. The
  mask-\textbf{definition} axis --- expert/sersic/perband/global
  reprojected to the IFU grid,
  \BreakablePath{run\_\discretionary{}{}{}sigma\_\discretionary{}{}{}e\_\discretionary{}{}{}mask\_\discretionary{}{}{}systematic.\discretionary{}{}{}py}
  --- gives ±4.58, which overlaps this; larger-of-two is kept, not
  added.)}
\item
  \textbf{Fit-window (±3.82):} 3 windows (wR3800\_5400 / wR4000\_5400 /
  w6500\_7500) × 3 SPS × N=500; peak-to-peak/2 (269.66 / 268.58 /
  276.22).
\item
  \textbf{Frame (±5.0):} structural --- vac/air per-SPS native-frame
  choice.
\item
  \textbf{Centering (±4.0):} 5 perturbed HST-mean centres (±0.4″ sweep;
  \BreakablePath{NOTES\_\discretionary{}{}{}centering\_\discretionary{}{}{}investigation}).
\item
  \textbf{Reduction-pass (±3.45):} half-Δ between the NEW and OLD
  cleaned cubes (269.62 vs 262.72); 2 reductions only --- refine if a
  3rd lands.
\item
  \textbf{R\_e-source (±5.13):} σ\_e re-measured over a 7-point R\_e
  grid bracketing the 2.097″ best-mask headline
  (\BreakablePath{run\_\discretionary{}{}{}sigma\_\discretionary{}{}{}e\_\discretionary{}{}{}Re\_\discretionary{}{}{}grid.\discretionary{}{}{}py}).
  \textbf{Adopted} = best-mask CoG family \{1.912/2.097/2.281″\}
  peak-to-peak/2 = (269.99−259.73)/2. The full grid incl.~CaHK+G 2.90″
  (±9.98) and CaHK+G's +12.4 deviation are cross-checks, NOT folded
  (CaHK+G is a different I-map definition, not an error on the
  photometric R\_e). Was ±6.13 at the old expert-mask headline.
\end{itemize}

Quote the budget table (§3.1) as the result; the iterative history that
\emph{led} to this masking + component set is the \textbf{internal audit
trail (M8→M12)} at the end of this document --- not for the manuscript
Methods section.

\hypertarget{reproducibility-infrastructure}{%
\subsubsection{§2.4.4 Reproducibility /
infrastructure}\label{reproducibility-infrastructure}}

\begin{itemize}
\tightlist
\item
  \textbf{Architecturally-independent estimators} (must not share
  intermediate state): \textbf{§6cum} cumulative I-weighted single-ppxf
  aperture (the Cappellari 2006 co-add-then-fit method = headline) ·
  \textbf{§7} discrete \textbf{Gültekin (2009) eq. 1} annular spaxel-sum
  (\BreakablePath{run\_\discretionary{}{}{}gultekin\_\discretionary{}{}{}mc},
  arc-filtered to R\textless R\_safe) · \textbf{nb07e} arc-spectrum
  subtraction. The three agree at \textless1σ (267.32 co-add-fit vs
  256.17 discrete, narrow).
\item
  \textbf{Caching + seed discipline.} Every expensive bootstrap writes a
  \BreakablePath{.\discretionary{}{}{}npz} keyed by its parameters
  (\BreakablePath{results/\discretionary{}{}{}<sweep>/\discretionary{}{}{}<params>\_\discretionary{}{}{}N\{N\}.\discretionary{}{}{}npz});
  a cached result is \textbf{never} re-run. Seeds are deterministic
  (§2.4.1 step 7) so re-runs bit-reproduce; \textbf{N=50 smoke before
  N=500 production}; joblib with BLAS=1/worker to avoid
  oversubscription.
\item
  \textbf{Deterministic values registry (single source of truth).}
  \BreakablePath{scripts/\discretionary{}{}{}paper\_\discretionary{}{}{}values.\discretionary{}{}{}py}
  recomputes \textbf{every} headline number from the result caches (with
  per-entry provenance + formula) → emits
  \BreakablePath{results/\discretionary{}{}{}PAPER\_\discretionary{}{}{}VALUES.\discretionary{}{}{}json}.
  The ApJL figures \textbf{load} that JSON (no hard-coded literals); the
  summary docs are validated by
  \BreakablePath{paper\_\discretionary{}{}{}values.\discretionary{}{}{}py -\discretionary{}{}{}-\discretionary{}{}{}check <files>},
  an anchored, precision-aware \textbf{drift linter} (parses only the
  canonical \BreakablePath{σ\_\discretionary{}{}{}e = N ± M} /
  \BreakablePath{log M⋆ = N} statements; dated snapshots carry a
  \BreakablePath{pv-\discretionary{}{}{}skip-\discretionary{}{}{}file}
  marker). Regenerate + verify:
  \BreakablePath{python scripts/\discretionary{}{}{}paper\_\discretionary{}{}{}values.\discretionary{}{}{}py \&\& python scripts/\discretionary{}{}{}paper\_\discretionary{}{}{}values.\discretionary{}{}{}py -\discretionary{}{}{}-\discretionary{}{}{}check *.\discretionary{}{}{}md}.
\end{itemize}

\begin{center}\rule{0.5\linewidth}{0.5pt}\end{center}

\hypertarget{results}{%
\section{§3 Results}\label{results}}

\hypertarget{figure-2-left-panel-the-integrated-spectrum-sed}{%
\subsection{Figure 2 (left panel) --- the integrated spectrum +
SED}\label{figure-2-left-panel-the-integrated-spectrum-sed}}

\begin{itemize}
\tightlist
\item
  \textbf{Left:} KCWI/KCRM Red-Low integrated, I(r)-convolved spectrum
  of the deflector elliptical, with the ppxf best fit overlaid (and
  inline residuals --- Rodrigo's standard nested-gridspec style). Loads
  from the wide-window fit + bootstrap pool.
\item
  \textbf{Right:} HST+JWST four-band SED + Bagpipes posterior model.
  Loads
  \BreakablePath{results/\discretionary{}{}{}bagpipes\_\discretionary{}{}{}sed\_\discretionary{}{}{}results.\discretionary{}{}{}npz}.
\item
  Figure file:
  \BreakablePath{AGEL0206\_\discretionary{}{}{}sigma\_\discretionary{}{}{}e\_\discretionary{}{}{}SED\_\discretionary{}{}{}final\_\discretionary{}{}{}wide.\discretionary{}{}{}pdf}
  (per outline) /
  \BreakablePath{results/\discretionary{}{}{}AGEL0206\_\discretionary{}{}{}spectra\_\discretionary{}{}{}SED\_\discretionary{}{}{}fit.\discretionary{}{}{}pdf}.
\end{itemize}

\hypertarget{velocity-dispersion-measurement}{%
\subsection{§3.1 Velocity Dispersion
Measurement}\label{velocity-dispersion-measurement}}

\begin{itemize}
\tightlist
\item
  \textbf{Final redshifts (to quote in Results).} The deflector systemic
  redshift from \textbf{Gaussian line fits} (6 absorption features,
  §2.2.2) is \textbf{z = 0.67564 ± 0.00033}, in agreement with the
  independent \textbf{ppxf stellar-kinematic} redshift \textbf{z =
  0.67555 ± 0.00002} (V\_sys = −16 ± 4 km/s relative to the line-fit
  value; headline R\_e=2.097″ aperture). \textbf{All three SPS libraries
  converge:} FSPS 0.67555 / EMILES 0.67553 / XSL 0.67556
  (\textasciitilde4 km/s spread) --- the signature of the air↔vacuum
  per-SPS frame fix (§2.4.1), which collapsed the legacy
  \textasciitilde110 km/s FSPS-vs-air offset. The 16 km/s line-fit↔ppxf
  offset is well within the ±59 km/s line-fit uncertainty (and the ppxf
  formal ±4 km/s carries an additional \textasciitilde±5 km/s per-SPS
  frame systematic). \textbf{We adopt z = 0.67564} (the line-fit value
  supersedes AGEL DR2 z = 0.67511). The lensed-\textbf{source} redshift
  is \textbf{z = 1.30263 ± 0.00003} (red-cube {[}O II{]} λλ3726,3729
  doublet). {[}§2.2.2;
  \BreakablePath{scripts/\discretionary{}{}{}redshift\_\discretionary{}{}{}verify.\discretionary{}{}{}py},
  \BreakablePath{run\_\discretionary{}{}{}sigma\_\discretionary{}{}{}e\_\discretionary{}{}{}Re\_\discretionary{}{}{}grid.\discretionary{}{}{}py}{]}
\item
  \textbf{Headline: σ\_e(\textless R\_e) = 267.31 ± 12.79 km/s} (asym
  −12.98/+12.61), the luminosity-weighted σ\_e within R\_e (Gültekin et
  al.~2009 eq. 1 definition) measured by the \textbf{Cappellari et
  al.~2006 (SAURON IV §2.3) co-add-then-fit method} --- a single pPXF
  fit to the I(r)-weighted R\textless R\_e aperture spectrum (R\_e =
  2.097″ best-mask CoG) --- plus the wild-bootstrap error pool.
  \textbf{Supersedes 269.62 ± 13.27 at R\_e=2.305″} --- the ``best mask
  throughout'' cascade (2026-06-11) adopts the best-mask R\_e=2.097″,
  lowering σ\_e by 2.3 km/s along the rising σ(R) profile (well within
  errors).
\item
  \textbf{Reproduce:} σ\_e at the headline aperture is
  \BreakablePath{results/\discretionary{}{}{}run\_\discretionary{}{}{}wide\_\discretionary{}{}{}sigma\_\discretionary{}{}{}e/\discretionary{}{}{}resys\_\discretionary{}{}{}best\_\discretionary{}{}{}mean/\discretionary{}{}{}}
  (R\_e=2.097″), produced by
  \BreakablePath{scripts/\discretionary{}{}{}run\_\discretionary{}{}{}sigma\_\discretionary{}{}{}e\_\discretionary{}{}{}Re\_\discretionary{}{}{}grid.\discretionary{}{}{}py -\discretionary{}{}{}-\discretionary{}{}{}n\_\discretionary{}{}{}bootstrap 500};
  the R\_e-source systematic (±5.13, best-mask CoG family) is from the
  same grid
  (\BreakablePath{results/\discretionary{}{}{}sigma\_\discretionary{}{}{}e\_\discretionary{}{}{}Re\_\discretionary{}{}{}grid\_\discretionary{}{}{}N500.\discretionary{}{}{}npz}).
\end{itemize}

\textbf{R\_e measurement and checks} (for the §3.1 paragraph): R\_e =
2.097″ = 14.76 kpc, the F140W (1.912″) + F200LP (2.281″)
\textbf{best-mask}-CoG mean (photutils-validated ±0.002″); supersedes
expert-mask 2.305″ and IFU-only 2.61″; masked CoG removes arc/companion
bias.

\textbf{σ\_e systematic budget} --- \emph{generated by
\BreakablePath{scripts/\discretionary{}{}{}paper\_\discretionary{}{}{}values.\discretionary{}{}{}py -\discretionary{}{}{}-\discretionary{}{}{}render};
do not hand-edit inside the markers.}

\begin{longtable}[]{@{}
  >{\raggedright\arraybackslash}p{(\columnwidth - 4\tabcolsep) * \real{0.33}}
  >{\raggedright\arraybackslash}p{(\columnwidth - 4\tabcolsep) * \real{0.33}}
  >{\raggedright\arraybackslash}p{(\columnwidth - 4\tabcolsep) * \real{0.33}}@{}}
\toprule
Component & ± km/s & Note \\
\midrule
\endhead
stat (N=500) & 4.51 & asym −5.04/+3.98; pooled 3 SPS × 15 deg
(marginalizes SPS + degree) \\
I-shape & 2.29 & 14-shape peak-to-peak/2 (266.79-271.37): the 10 raw
maps + 4 arc-free PSF-matched (Sérsic/filled→1.27" conv); was ±2.27 over
10 \\
F200 mask & 6.65 & (w00-w100)/2; larger-of-two vs mask-approach 5.85 (no
double-count) \\
frame (vac/air) & 5.00 & carried constant \\
centering & 4.00 & carried constant \\
fit-window & 3.82 & peak-to-peak/2 across 3 fit windows \\
reduction-pass & 3.45 & half-Δ between NEW and OLD reductions \\
R\_e-source (best-mask grid) & 5.13 & peak-to-peak/2 across best-mask
CoG family \{F140W 1.91``, mean 2.10'', F200LP 2.28``\} at the 2.097''
headline (user 2026-06-11; CaHK+G \& full grid are cross-checks, not
folded) \\
\textbf{TOTAL (sym)} & \textbf{12.79} & quadrature (sys ±11.97 ⊕
stat) \\
\textbf{TOTAL (asym)} & \textbf{−12.98 / +12.61} & preserves stat-side
skew \\
\bottomrule
\end{longtable}

Cross-checks (not added): arc-mask-definition ±4.58 (overlaps F200 mask)
· full-grid R\_e-source ±9.98 (incl.~CaHK+G, conservative ceiling) ·
CaHK+G(2.90″) deviation +12.38 km/s.

\emph{(Per-component sweep detail --- the per-window / per-estimator
values behind each ± --- is in §2.4.3.)}

\begin{itemize}
\tightlist
\item
  \textbf{No separate SPS line} is folded in: the SPS-library spread
  \textbf{collapsed from \textasciitilde26 km/s at the narrow window}
  (FSPS 253.0 \textless{} EMILES 267.5 \textless{} XSL 279.7) \textbf{to
  \textasciitilde4 km/s at the wide window} (FSPS 271.3 / EMILES 267.2 /
  XSL 270.4;
  \BreakablePath{PAPER\_\discretionary{}{}{}VALUES.\discretionary{}{}{}json:per\_\discretionary{}{}{}sps\_\discretionary{}{}{}p50})
  --- a key argument for the wide window; it now sits inside the pooled
  stat width (between-SPS ±2.04).
\item
  Quadrature is deliberately conservative (components not strictly
  independent).
\item
  \textbf{R\_e-source systematic IS now folded in (±6.13, M12
  2026-06-08):} re-measured at the wide window (nb15) and included in
  the budget above --- the narrow-window ±16.9 is superseded. The full
  4-estimator spread was adopted (user decision); the CaHK+G 2.90″
  estimator drives it (+10 km/s = rising σ(R), see below), while the
  light-R\_e family alone gives ±2.50.
\end{itemize}

\textbf{Systematics paragraph order} (per outline): first R\_e and I(r)
(I-shape ±2.27, R\_e-source flagged), then the spectrum (mask ±6.65,
frame ±5.0, window ±3.82, reduction ±3.45, centering ±4.0).

\textbf{Cross-checks} (narrow window, architecturally independent,
superseded for headline but valid):

\begin{itemize}
\tightlist
\item
  §6cum cumulative I-weighted aperture ppxf: 267.32 ± 24;
  σ\_e(\textless R\_e/2) ≈ 225.78; σ\_e(\textless R\_e/8) ≈ 209.18.
\item
  §7 discrete Gültekin annular sum (equal-N inside R\_safe=3R\_e/4):
  256.17 −13.0/+12.7.
\item
  §7b flat-σ extrapolation: 274.37 −16.2/+17.4.
\item
  nb07e arc-spectrum subtraction matches §6cum bit-identically
  (\textless1 km/s).
\item
  Narrow Ca H+K + G window (nb09d): 267.95 ± 30.10 --- headline within
  \textasciitilde2 km/s.
\end{itemize}

\textbf{Headline evolution} (mask-audit arc): 268.98 (clean) → 271.87
(+He I, M8) → 269.62 (+M10 sky audit) → ±11.77 (M11 re-derivation) →
\textbf{±13.27 (M12, +R\_e-source ±6.13, 2026-06-08)}. He I pushes σ up
\textasciitilde2.9, the 9 sky bands pull it down \textasciitilde2.2
(near-cancel); central value 269.62 stable across M10→M12 (only the
error grew).

\textbf{Compare with Greene (2016/2020) and KH13 choices:}
σ\_e(\textless R\_e) here uses the luminosity-weighted single-aperture
σ\_e (Cappellari et al.~2006 SAURON IV §2.3 co-add-then-fit; Gültekin
2009 eq. 1 definition), exactly as adopted by Kormendy \& Ho 2013 (M•--σ
eq. 3) and Greene et al.~2020 ARA\&A (fig.~5). We dropped
R\textless R\_e/8 (JFK95 σ\_c) because seeing under-resolves it. For a
pure elliptical, galaxy R\_e ≈ bulge R\_e, so KH13's ``bulge R\_e''
equals the measured total R\_e --- no bulge/disk decomposition needed.

\textbf{σ\_e physical-meaning discussion:} the deflector is a
dispersion-supported \textbf{elliptical bulge}; σ(r) should
\emph{decrease} gently with radius --- any rising σ(r) in raw output is
a continuum-dilution artifact from the z=1.302 arc, not a physical
gradient. Resolved kinematics (nb11): 17 central spaxels at S/N≥5, σ ∈
{[}144, 251{]} km/s (median 201), declining with R, no coherent rotation
above noise.

\hypertarget{stellar-mass-estimate}{%
\subsection{§3.2 Stellar Mass Estimate}\label{stellar-mass-estimate}}

\begin{itemize}
\tightlist
\item
  \textbf{HEADLINE: log(M⋆/M☉) = 11.50 +0.09/−0.14} (aperture-corrected
  total, \textbf{matched 2 R\_e aperture}, 10\% flux floor; Planck 2015;
  \BreakablePath{scripts/\discretionary{}{}{}aperture\_\discretionary{}{}{}matched\_\discretionary{}{}{}photometry.\discretionary{}{}{}py},
  \BreakablePath{results/\discretionary{}{}{}aperture\_\discretionary{}{}{}matched\_\discretionary{}{}{}photometry.\discretionary{}{}{}npz}).
  Total-light mass: empirical flux in a \textbf{matched 2 R\_e
  elliptical aperture} (same physical region all 4 bands; single-Sérsic
  deflector model; color/morph-gated WCS-registered arc mask + companion
  mask), corrected to total by \textbf{adding the Sérsic model's
  beyond-aperture light} (the GAMA/Taylor+2011 \BreakablePath{fluxscale}
  approach; additive variant). Methodologically: \textbf{Taylor et
  al.~2011} (GAMA aperture→total for M⋆), \textbf{Sonnenfeld et
  al.~2013} (SL2S lens-deflector Sérsic photometry, mask arc +
  interlopers), \textbf{Graham \& Driver 2005} (Sérsic total formalism).
\item
  \textbf{Report FIVE estimators (empirical → model), M⋆ for each, at 2
  R\_e:} \textbar{} estimator \textbar{} log M⋆ \textbar{} what it adds
  \textbar{}
  \textbar-----------\textbar--------\textbar--------------\textbar{}
  \textbar{} raw (empirical aperture, masked) \textbar{} \textbf{11.21}
  \textbar{} nothing --- lower bound \textbar{} \textbar{} raw + apcorr
  (most-empirical total) \textbar{} \textbf{11.38} \textbar{} model
  \textbf{wings} only (masked interior NOT filled) \textbar{} \textbar{}
  filled (within aperture) \textbar{} \textbf{11.38} \textbar{} model
  \textbf{fill} of masked pixels \textbar{} \textbar{} \textbf{total
  (aperture-corrected) --- headline} \textbar{} \textbf{11.50}
  \textbar{} fill + wings \textbar{} \textbar{} Sérsic-total (pure
  model) \textbar{} \textbf{11.43} \textbar{} everything modeled
  \textbar{} (all five on Planck 2015 --- the +0.0282 dex D\_L² rescale
  vs the prior H₀=70 values 11.22/11.35/11.36/11.47/11.41.) The
  \textbf{total converges across 1/2/2.5 R\_e} (11.51/11.50/11.49) → the
  aperture correction is internally consistent. raw converges 2↔2.5 R\_e
  (11.21/11.23).
\end{itemize}

\textbf{M⋆ error budget} --- \emph{generated by
\BreakablePath{scripts/\discretionary{}{}{}paper\_\discretionary{}{}{}values.\discretionary{}{}{}py -\discretionary{}{}{}-\discretionary{}{}{}render};
do not hand-edit inside the markers.} Unlike σ\_e (a single
named-component quadrature), M⋆ is reported in \textbf{two framings}:
the headline error carried on the adopted total-light estimator, and the
Sérsic-only named-component quadrature kept as a model-systematic
cross-check (so model choices are not double-counted into the
empirical-aperture mass).

\textbf{(1) Headline --- log(M⋆/M☉) = 11.50 (aperture-corrected total,
matched 2 R\_e, 10\% floor):}

\begin{longtable}[]{@{}
  >{\raggedright\arraybackslash}p{(\columnwidth - 4\tabcolsep) * \real{0.33}}
  >{\raggedright\arraybackslash}p{(\columnwidth - 4\tabcolsep) * \real{0.33}}
  >{\raggedright\arraybackslash}p{(\columnwidth - 4\tabcolsep) * \real{0.33}}@{}}
\toprule
Component & ± dex & Note \\
\midrule
\endhead
stat (Bagpipes posterior) & +0.094 / −0.141 & asymmetric; age--dust--M/L
outshining low-tail (present in all 5 estimators incl.~raw) \\
masking-approach & 0.086 & under-arc (raw↔filled) 0.057 ⊕ mask-def
(per-band↔global) 0.028 \\
\textbf{REPORTED} & \textbf{+0.09 / −0.14} & the posterior is quoted as
the headline error; the ±0.09 dex masking systematic is reported
alongside \\
\bottomrule
\end{longtable}

\textbf{(2) Sérsic-only estimator (log(M⋆/M☉) = 11.43) --- model-path
systematic envelope (cross-check), the named-component analogue of the
σ\_e budget:}

\begin{longtable}[]{@{}lll@{}}
\toprule
Component & ± dex & Note \\
\midrule
\endhead
stat & 0.121 & Bagpipes posterior width \\
mask & 0.125 & arc-mask choice (expert↔global Sérsic-total) ---
dominant \\
model-form & 0.057 & single-Sérsic vs alternative model form \\
flux-floor & 0.034 & 10\%↔20\% flux floor \\
fit-param & 0.029 & Sérsic-n / fit-parameter spread \\
apcorr-recon & 0.026 & apcorr ↔ pure-model reconstruction \\
\textbf{TOTAL} & \textbf{0.19} & stat ±0.12 ⊕ sys ±0.15
(mask-dominated) \\
\bottomrule
\end{longtable}

Five M⋆ estimators (empirical→model, matched 2 R\_e, Planck 2015): raw
\textbf{11.21} (empirical lower bound) · raw+apcorr \textbf{11.38} ·
filled \textbf{11.38} · \textbf{total 11.50 (headline)} · Sérsic-total
\textbf{11.43}.

\emph{(Per-component provenance --- which cache + sweep produced each ±
--- is in the bullets below and in §2.1.1b / §2.4.3. The σ\_e budget is
in §3.1.)}

\begin{itemize}
\tightlist
\item
  \textbf{Sérsic-only (full-model) systematic budget = ±0.19 dex}
  (\BreakablePath{scripts/\discretionary{}{}{}sersic\_\discretionary{}{}{}total\_\discretionary{}{}{}systematic.\discretionary{}{}{}py};
  elliptical multi-start fit, 2026-06-11): stat ±0.121 ⊕ sys ±0.147,
  with \textbf{arc-mask choice ±0.125 dominant}, model-form ±0.057,
  flux-floor ±0.034, Sérsic-fit-n ±0.029, apcorr↔pure-model ±0.026. (Was
  ±0.17 with the circular fit; the mask term grew because the elliptical
  model widens the expert↔global Sérsic-total gap.)
\item
  \textbf{Apcorr chain validated vs established codes (2026-06-11,
  \BreakablePath{scripts/\discretionary{}{}{}validate\_\discretionary{}{}{}apcorr\_\discretionary{}{}{}established.\discretionary{}{}{}py};
  petrofit/statmorph absent → photutils + scipy incomplete-gamma
  references):} (1) b\_n Ciotti\&Bertin99 vs exact
  \BreakablePath{gammaincinv(2n,½)} ≤0.05\%; (2) the Sérsic total-flux
  formula (Graham\&Driver05 eq.4) vs a numeric render-to-20 R\_e
  ≤0.03\%; (3) the aperture correction --- rendered-model enclosed light
  vs the \textbf{analytic Sérsic incomplete-gamma law}
  \BreakablePath{γ(2n, b\_\discretionary{}{}{}n·2\^\{1/\discretionary{}{}{}n\})/\discretionary{}{}{}Γ(2n)}
  (Graham\&Driver05) --- Δ≤0.0007;

  \begin{enumerate}
  \def\labelenumi{(\arabic{enumi})}
  \setcounter{enumi}{3}
  \tightlist
  \item
    the empirical hard-edge aperture sum
    (\BreakablePath{in\_\discretionary{}{}{}aperture}) vs photutils
    \BreakablePath{EllipticalAperture(method='exact')} ≤0.09\%; (5) the
    finite-cutout truncation of
    \BreakablePath{sum(model\_\discretionary{}{}{}full)} (the apcorr
    denominator) vs the to-∞ total ≤0.19\% (≤0.002 mag). \textbf{No
    production bug; M⋆ unchanged.}
  \end{enumerate}
\item
  \textbf{Single-Sérsic ellipticity fix (2026-06-11).} The deflector
  light model is a single Sérsic per band fit on the best mask. The
  fitter originally had a single low-ellipticity start that railed into
  a spurious \textbf{circular} χ² minimum (b/a=1 for all bands),
  under-reporting the shape and biasing the \BreakablePath{(1−ellip)}
  total flux. Fixed with a \textbf{multi-start} (flux-weighted-moment
  seed à la SExtractor/GALFIT

  \begin{itemize}
  \tightlist
  \item
    a PA grid + circular fallback, lowest-residual). It recovers the
    true shape; \textbf{M⋆ headline (total) is unchanged at 11.47}
    (empirical-aperture-dominated). Validated by injection-recovery
    against astropy's peer-reviewed \BreakablePath{Sersic2D}
    (\BreakablePath{scripts/\discretionary{}{}{}validate\_\discretionary{}{}{}sersic\_\discretionary{}{}{}fitter\_\discretionary{}{}{}synthetic.\discretionary{}{}{}py}):
    clean recovery for the deflector regime (n≈1.2--1.6, b/a≤0.85), with
    a realistic \textbf{b/a uncertainty \textasciitilde±0.06} (the
    formal bootstrap is too tight). GALFIT/imfit/petrofit/statmorph are
    not installed (conda-env policy).
  \end{itemize}
\item
  \textbf{The M⋆ posterior has a longer LOW tail} (total 11.50
  +0.09/−0.14). This is \textbf{NOT} the total-Sérsic estimator --- it
  is present in all five estimators including the purely empirical
  \BreakablePath{raw}. It is the \textbf{age--dust--M/L ``outshining''
  degeneracy} in the 4-band Bagpipes fit: the low-M⋆ tail correlates
  with young age (r=+0.90), high dust (r=−0.58), and high sSFR (r=−0.83)
  --- young+dusty low-M/L solutions reproduce the same broadband points.
  We keep the exponential-SFH prior and report the asymmetric posterior
  as-is (the tail is real model uncertainty, not an estimator artifact).
\item
  \textbf{Per-band single-Sérsic parameters (appendix table;
  \BreakablePath{scripts/\discretionary{}{}{}sersic\_\discretionary{}{}{}parameter\_\discretionary{}{}{}table.\discretionary{}{}{}py}):}
\end{itemize}

\begin{longtable}[]{@{}
  >{\raggedright\arraybackslash}p{(\columnwidth - 16\tabcolsep) * \real{0.11}}
  >{\raggedright\arraybackslash}p{(\columnwidth - 16\tabcolsep) * \real{0.11}}
  >{\raggedright\arraybackslash}p{(\columnwidth - 16\tabcolsep) * \real{0.11}}
  >{\raggedright\arraybackslash}p{(\columnwidth - 16\tabcolsep) * \real{0.11}}
  >{\raggedright\arraybackslash}p{(\columnwidth - 16\tabcolsep) * \real{0.11}}
  >{\raggedright\arraybackslash}p{(\columnwidth - 16\tabcolsep) * \real{0.11}}
  >{\raggedright\arraybackslash}p{(\columnwidth - 16\tabcolsep) * \real{0.11}}
  >{\raggedright\arraybackslash}p{(\columnwidth - 16\tabcolsep) * \real{0.11}}
  >{\raggedright\arraybackslash}p{(\columnwidth - 16\tabcolsep) * \real{0.11}}@{}}
\toprule
Band & \(\lambda_{\rm piv}\) (Å) & \(r_e\) (″) & \(r_e\) (kpc) & \(n\) &
\(b/a\) & PA (°) & \(\mu_e\) (mag/″²) & \(m_{\rm AB}^{\rm tot}\) \\
\midrule
\endhead
F200LP & 4972 & 2.002±0.122 & 14.09±0.86 & 1.24±0.15 & ≳0.95† & ---† &
25.23 & 20.92±0.07 \\
F140W & 13923 & 1.765±0.015 & 12.42±0.11 & 1.44±0.15 & 0.86±0.06 & 4±20
& 22.36 & 18.42±0.01 \\
F150W2 & 16720 & 1.643±0.005 & 11.57±0.03 & 1.22±0.15 & 0.80±0.06 &
11±20 & 21.91 & 18.29±0.00 \\
F322W2 & 32470 & 1.757±0.010 & 12.37±0.07 & 1.59±0.15 & 0.85±0.06 & 6±20
& 21.79 & 17.82±0.01 \\
\bottomrule
\end{longtable}

† F200LP: the deflector is faint in this band, so the single-Sérsic
ellipticity/PA are unconstrained (consistent with circular); the IR
bands set the shape. Errors are bootstrap, with \(b/a\)/PA/\(n\) floored
by the injection-recovery scatter; the mask-choice systematic (§2.1.1b)
dominates and is separate. The deflector is \textbf{mildly elliptical
(b/a≈0.80--0.86, PA≈4--11° E of N)} with low Sérsic index (n≈1.2--1.6)
--- consistent with the b/a=0.75 adopted for the aperture from the
independent light fit. - \textbf{Aperture provenance / why this
supersedes 11.36:} the old per-band apertures were hand-drawn in the GUI
and \textbf{inconsistent} --- F200LP a=1.19″ (\textasciitilde0.5 R\_e)
captured only \textbf{17\%} of the light vs IR \textasciitilde46\% ---
so F200LP was under-measured → SED artificially red → M⋆ biased
\textbf{HIGH} (11.36). The matched aperture measures the same region in
every band → bluer SED → lower M/L → the empirical raw drops to 11.22. -
\textbf{Flux uncertainties:} the \emph{propagated statistical} errors
are \textbf{0.03--1.6\%} (F322W2 0.03\%, F150W2 0.05\%, F140W 0.4\%,
F200LP 1.6\%) --- 6--362× below the adopted \textbf{10\% floor}. The
10\% is a deliberate \textbf{systematic floor} (zeropoint, aperture,
drizzle correlated noise \textasciitilde1.3--2×, SED-model mismatch),
not the measurement error --- report the true errors and state the floor
(\textasciitilde5\% is better-motivated; precedent in GAMA/Taylor+2011,
whose error budgets are calibration-dominated, not photon-noise).
\textbf{Floor sensitivity
(\BreakablePath{results/\discretionary{}{}{}aperture\_\discretionary{}{}{}floor5\_\discretionary{}{}{}check.\discretionary{}{}{}npz}):}
5\%↔10\% shifts every estimator's central by ≤+0.02 dex (headline total
10\%=11.50, 5\%=11.52, Planck 2015) with tighter stat at 5\% --- well
within the ±0.12 systematic, so the headline is robust to the floor;
10\% kept for the headline. - \textbf{Superseded values (cross-checks):}
mismatched-aperture 11.36 (F200LP-biased high); 2-component-mask 11.16
+0.31 (over-masked → biased low); expert-aperture 11.33. -
\textbf{``Typical elliptical at z\textasciitilde0.7'':} M⋆ ≈ 1.6 × 10¹¹
M☉ (empirical raw) to ≈ 3.2 × 10¹¹ (aperture-corrected total headline)
at z = 0.6756, passive/quiescent SFH (mass-weighted age \textasciitilde5
Gyr, low sSFR). Place on the z≈0.5--1.0 mass-size relation (van der Wel
et al. 2014; Mowla et al.~2019) and the strong-lens-deflector analog
sample (Sonnenfeld et al.~2015, SL2S V --- best z-match). {[}notebook 10
verified citations{]} - \textbf{X-ray / quiescence:} \textbf{GAP G2 ---
no X-ray data in repo.} Source externally if you want the ``truly
quiescent'' claim; otherwise frame from the passive SED only.

\hypertarget{mass-estimates-from-lens-modeling}{%
\subsection{§3.3 Mass Estimates from Lens
Modeling}\label{mass-estimates-from-lens-modeling}}

\begin{quote}
\textbf{G1 --- sourced from the Ferrami et al.~draft (2026-06-12); see
§2.3 for method + caveats.} \textbf{Every number below is PROVISIONAL}
--- the draft stamps the evidence table, the M•--σ figure, and the
potential-correction figure
\BreakablePath{PLACEHOLDER UNTIL FINAL RUNS ARE COMPLETED}. The draft is
also internally inconsistent on the headline BH mass (several values
coexist, see flag) because the final sampling has not landed. \textbf{Do
not freeze any lens-model number into our draft until Ferrami's final
runs complete.} All masses on Ferrami's \textbf{Planck-2015 cosmology
(H₀=67.7, Ω\_m=0.302)} --- convert to our (H₀=70, Ω\_m=0.3) frame before
any joint M•--σ / M•--M⋆ plot (§2.3 reconciliation flags).
\end{quote}

\hypertarget{headline-result}{%
\subsubsection{§3.3.1 Headline result}\label{headline-result}}

\begin{itemize}
\tightlist
\item
  \textbf{Canonical BH mass = the combined free-BH posterior:}
  \textbf{M\_• = 5.2 ₋₁.₃/+₁.₅ × 10⁸ M☉} (≈ \textbf{5.2 ± 1.4 × 10⁸
  M☉}). This is the model-averaged posterior across the three
  \textbf{free-BH} models (BH position free within 0.5″ of the light
  centre), which are the \textbf{highest-evidence family} in Table 1 ---
  cEPL+free BH (Δln𝓔 = 0), bEPL+free (−0.08), sEPL+free (−1.43).
  Per-model medians: cEPL \textbf{5.51 ± 1.40}, sEPL \textbf{5.01
  ₋₁.₂₀/+₁.₅₀}, bEPL \textbf{5.11 ₋₁.₃₀/+₁.₅₀} (×10⁸ M☉); they agree to
  \textasciitilde10 \%, so equal-weight and evidence-weighted averaging
  give the same answer (5.22 vs 5.29 ×10⁸). \textbf{Conservative 2σ
  envelope across the three: 2.6--8.3 × 10⁸ M☉} (the draft's own
  ``min/max of the 2σ percentiles'' prescription). Uniform prior
  \textbf{{[}10⁸, 10¹¹{]} M☉}. \emph{{[}User decision 2026-06-12: adopt
  the free-BH combined posterior, NOT the fixed-BH
  \textasciitilde2×10⁹.{]}}

  \begin{itemize}
  \tightlist
  \item
    \textbf{Provenance / caveat:} combined here from the Ferrami
    \textbf{Table 1} summary medians ± errors as a split-normal mixture
    (\BreakablePath{scripts/\discretionary{}{}{}bh\_\discretionary{}{}{}mass\_\discretionary{}{}{}combine.\discretionary{}{}{}py},
    seeded; equal-weight median 5.22, evidence-weighted 5.29 ×10⁸). The
    \textbf{rigorous} combination stacks Giovanni's actual
    \textbf{MultiNest chains} --- flagged as a TODO, but the three
    free-BH posteriors are near-identical so the summary-level result is
    robust.
  \end{itemize}
\item
  \textbf{Detection significance \textgreater{} 4σ:} models with a
  central \textbf{point mass} have higher Bayesian evidence than
  smooth-only (no-BH) models at \textbf{\textgreater{} 4σ} --- the
  radial image \emph{requires} a central point-like mass. (The no-BH and
  extended-mass SIS/NFW perturbers are all disfavoured by Δln𝓔 ≈ −4 to
  −6; §3.3.3.)
\item
  \textbf{⚠ Supersedes the earlier ``\textasciitilde2×10⁹ fixed-BH''
  framing.} The Ferrami draft \emph{prose} (intro/results + best-model
  figure caption) is written around the \textbf{fixed}-BH solutions
  (\textasciitilde1.6--2.62×10⁹ M☉), but those have \textbf{lower
  evidence} than the free-BH runs. We adopt the free-BH combined
  posterior (5.2×10⁸); the fixed-BH values are kept below only as the
  lower-evidence alternative. \textbf{Giovanni's prose still calls
  sEPL+fixed-PM the ``best model'' --- this must be reconciled} (the
  data, when the BH is free to move, prefer a \textasciitilde3--4× lower
  mass).
\end{itemize}

\hypertarget{best-model-why-the-slope-matters}{%
\subsubsection{§3.3.2 Best model + why the slope
matters}\label{best-model-why-the-slope-matters}}

\begin{itemize}
\tightlist
\item
  \textbf{Highest-evidence model = a FREE point mass on the cEPL (or
  bEPL) main halo} --- Δln𝓔 = 0 (cEPL+free) / −0.08 (bEPL+free) in Table
  1, BH free within 0.5″ of the light centre → our canonical M\_• =
  5.2×10⁸ (§3.3.1). \textbf{⚠ The draft prose instead designates sEPL +
  a FIXED point mass as the ``best model''} (and the best-model figure
  caption shows 2.62×10⁹) --- a \textbf{fixed-vs-free inconsistency}:
  the fixed runs have \emph{lower} evidence (Δln𝓔 ≈ −3.0 to −3.65) yet
  are what the prose/figure quote. We adopt the free-BH result; the
  final runs / Giovanni's prose must resolve this. (Evidence is quoted
  as Δln𝓔; the placeholder table is not renormalised to any single
  reference.)
\item
  \textbf{The main halo is markedly sub-isothermal: γ ≈ 1.31 ± 0.08}
  (θ\_E ≈ 1.36″). A shallow total slope is \emph{what produces the
  bright radial arc extending to the deflector centre} --- the lens
  environment (it sits \textasciitilde100 kpc from the cluster
  \textbf{ACT-CL J0206.2−0114}) flattens the projected mass. This is the
  ``shape of the mass profile matters'' argument: the radial image only
  stays bright down to the centre for a shallow slope, and it is the
  radial image that the BH perturbs. {[}Ferrami draft §Data, §Mass model
  comparison{]}
\end{itemize}

\hypertarget{model-comparison-table-15-models-placeholder-ferrami-draft-table-1}{%
\subsubsection{§3.3.3 Model-comparison table (15 models --- PLACEHOLDER,
Ferrami draft Table
1)}\label{model-comparison-table-15-models-placeholder-ferrami-draft-table-1}}

3 smooth main-halo families (sEPL/cEPL/bEPL) × 5 perturbers (free BH /
fixed BH / free NFW / free SIS / no BH). Δln 𝓔 are as printed in the
draft (NOT yet renormalised to the sEPL+fixed-PM reference):

\begin{longtable}[]{@{}
  >{\raggedright\arraybackslash}p{(\columnwidth - 6\tabcolsep) * \real{0.25}}
  >{\raggedright\arraybackslash}p{(\columnwidth - 6\tabcolsep) * \real{0.25}}
  >{\raggedright\arraybackslash}p{(\columnwidth - 6\tabcolsep) * \real{0.25}}
  >{\raggedright\arraybackslash}p{(\columnwidth - 6\tabcolsep) * \real{0.25}}@{}}
\toprule
Smooth model & Δln 𝓔 & d.o.f. & Perturber posterior \\
\midrule
\endhead
\textbf{cEPL + free BH} \emph{(highest evidence → canonical family)} &
\textbf{0} & 37 & \textbf{M\_• = 5.51 ± 1.40 ×10⁸ M☉} \\
cEPL + fixed BH & −3.25 & 35 & M\_• = 1.61 ± 0.21 ×10⁹ M☉ \\
cEPL + free NFW & −4.88 & 38 & M\_vir = 6.48 ×10⁹, c\_vir = 40 \\
cEPL + free SIS & −5.26 & 38 & θ\_E = 0.041″, γ = 2.59 \\
cEPL (no BH) & −5.06 & 34 & --- \\
\textbf{sEPL + free BH} \emph{(canonical family)} & −1.43 & 36 &
\textbf{M\_• = 5.01 ₋₁.₂₀/+₁.₅₀ ×10⁸ M☉} \\
sEPL + fixed BH \emph{(draft prose ``best model'' --- lower evidence)} &
−3.65 & 34 & M\_• = 2.10 ⁺⁰·¹⁹₋₀.₃₀ ×10⁹ M☉ \\
sEPL + free NFW & −4.76 & 37 & M\_vir = 2.65 ×10⁹, c = 91 \\
sEPL + free SIS & −5.80 & 37 & θ\_E = 0.027″, γ = 2.47 \\
sEPL (no BH) & −4.78 & 33 & --- \\
\textbf{bEPL + free BH} \emph{(canonical family)} & −0.08 & 38 &
\textbf{M\_• = 5.11 ₋₁.₃₀/+₁.₅₀ ×10⁸ M☉} \\
bEPL + fixed BH & −3.04 & 36 & M\_• = 1.86 ×10⁹ M☉ \\
bEPL + free NFW & −4.35 & 39 & M\_vir = 5.78 ×10⁹, c = 127 \\
bEPL + free SIS & −4.82 & 39 & θ\_E = 0.025″, γ = 2.46 \\
bEPL (no BH) & −5.14 & 35 & --- \\
\bottomrule
\end{longtable}

\begin{itemize}
\tightlist
\item
  \textbf{Pattern (robust across all 3 halo families):} the
  \textbf{no-BH} and \textbf{free-SIS/free-NFW} (extended central mass)
  models are all disfavoured by Δln𝓔 ≈ −4 to −6; a \textbf{point mass}
  is preferred. \textbf{The FREE-BH runs carry the HIGHEST evidence}
  (Δln𝓔 0 / −0.08 / −1.43) and cluster at \textbf{\textasciitilde5×10⁸
  M☉ → the canonical combined value (§3.3.1)}; the \emph{fixed}-BH runs
  have \textbf{lower} evidence (Δln𝓔 −3.0 to −3.65) and land
  \textasciitilde3--4× higher (1.6--2.1×10⁹). Earlier drafts of this
  fact sheet treated the fixed value as headline --- \textbf{corrected
  2026-06-12 to the free-BH combined posterior.}
\end{itemize}

\hypertarget{best-model-parameters-ferrami-draft-table-2-placeholder-headed-cepl-fixed-bh-run}{%
\subsubsection{\texorpdfstring{§3.3.4 Best-model parameters (Ferrami
draft Table 2 --- PLACEHOLDER; headed ``cEPL'', \textbf{FIXED-BH
run})}{§3.3.4 Best-model parameters (Ferrami draft Table 2 --- PLACEHOLDER; headed ``cEPL'', FIXED-BH run)}}\label{best-model-parameters-ferrami-draft-table-2-placeholder-headed-cepl-fixed-bh-run}}

\begin{quote}
\textbf{Note:} this parameter table is a \textbf{fixed-BH} solution
(M\_• = 1.61×10⁹, the lower-evidence run); the \textbf{canonical} BH
mass is the \textbf{free-BH combined posterior 5.2×10⁸} (§3.3.1). The
smooth-halo parameters below (θ\_E, γ, shear, light) are essentially
common to the free-BH solution --- only the point-mass row differs.
\end{quote}

\begin{longtable}[]{@{}
  >{\raggedright\arraybackslash}p{(\columnwidth - 4\tabcolsep) * \real{0.33}}
  >{\raggedright\arraybackslash}p{(\columnwidth - 4\tabcolsep) * \real{0.33}}
  >{\raggedright\arraybackslash}p{(\columnwidth - 4\tabcolsep) * \real{0.33}}@{}}
\toprule
Component & Parameter & Posterior \\
\midrule
\endhead
Main deflector & θ\_E & 1.360 ⁺⁰·¹⁰₋₀.₀₈ ″ \\
& γ (slope) & \textbf{1.31 ± 0.08} (sub-isothermal) \\
& r\_c (core) & 0.0012″ \\
& f (axis ratio) & 0.826 \\
& φ & 10.8° \\
Secondary deflector (SIS) & θ\_E & 0.194″ \\
Tertiary deflector (SIS) & θ\_E & 0.017″ \\
External shear & Γ & 0.095 ± 0.003 \\
& θ\_Γ & −14.8° \\
\textbf{Point mass} \emph{(fixed-BH; NOT canonical)} & \textbf{M\_•} &
1.61 ⁺⁰·²¹₋₀.₂₁ ×10⁹ M☉ → \textbf{use free-BH 5.2×10⁸ (§3.3.1)} \\
Lens light (Sérsic) & R\_eff & 0.051″ \\
& n\_s & 1.372 \\
& f (axis ratio) & 0.834 \\
Source light (Delaunay) & λ & 46798 \\
& w\_a, w\_b & 1.728, 0.451 \\
\bottomrule
\end{longtable}

(Circular multipole and lens-light multipole coefficients a\_m/b\_m,
m=1,3,4 are all \textbar a\textbar,\textbar b\textbar{} ≲ 0.012 with
Gaussian{[}0,0.01{]} priors --- see draft Table 2 for the full list.)

\hypertarget{mux3c3-placement-and-prior-radial-arc-bh-lenses}{%
\subsubsection{§3.3.5 M•--σ placement and prior radial-arc BH
lenses}\label{mux3c3-placement-and-prior-radial-arc-bh-lenses}}

\begin{itemize}
\tightlist
\item
  The draft places M\_• against the \textbf{local M•--σ relation of van
  den Bosch (2016, M\_σ FP)}, with the only two prior \textbf{radial-arc
  / radial-image SMBH lens} measurements overplotted:
  \textbf{Nightingale et al. (z = 0.17)} and \textbf{Melo-Carneiro et
  al.~(z = 0.45)} --- both \textasciitilde2σ \emph{above} the local
  relation. \textbf{DES J0206−0114 (z = 0.675) is the highest-redshift
  and most centrally-extended} of the three.
\item
  \textbf{⚠ With the canonical M\_• = 5.2×10⁸ (free-BH), the placement
  SENSE flips vs the old \textasciitilde2×10⁹.} Rough \textbf{KH13
  cross-check} (log M• = 8.49 + 4.38·log(σ\_e/200), intrinsic scatter
  0.29 dex): at σ\_e = 267 km/s KH13 predicts M• ≈ \textbf{1.1×10⁹}, so
  our 5.2×10⁸ (log M• = 8.72) sits \textbf{≈0.3 dex (\textasciitilde1σ)
  BELOW} the local relation --- i.e.~\emph{slightly under-massive} for
  its σ, the \textbf{opposite} of the fixed-BH \textasciitilde2×10⁹
  (which would sit \emph{above}, like Nightingale/Melo-Carneiro).
  \emph{(This is a labeled rough calc on KH13, NOT the paper result; the
  actual placement vs vdB16 is Giovanni's figure --- the narrative
  ``above the relation'' framing in the draft prose needs reconciling
  with the free-BH value.)}
\item
  The M•--σ figure is a \textbf{PLACEHOLDER}; the σ comes from
  \textbf{our companion paper} (Córdova Rosado et al., in prep. --- σ\_e
  = \textbf{267 ± 13 km/s}, this repo). The draft's conservative SMBH
  estimate is the min/max of the \textbf{2σ percentiles} across the
  perturber posteriors (= our \textbf{2.6--8.3×10⁸} envelope, §3.3.1).
  \textbf{This is exactly the 4D (M•, σ, M⋆, z) joint constraint
  PROJECT\_BRIEF.md flags} --- and now that both papers share the
  Planck-2015 cosmology (§2.3), the join is direct.
\end{itemize}

\hypertarget{quiescence-evidence-also-fills-gap-g2-x-ray-radio}{%
\subsubsection{§3.3.6 Quiescence evidence (also fills gap G2 --- X-ray /
radio)}\label{quiescence-evidence-also-fills-gap-g2-x-ray-radio}}

The Ferrami draft §Data carries the multiwavelength quiescence argument
that was gap \textbf{G2} in this repo (no longer external-only):

\begin{itemize}
\tightlist
\item
  \textbf{No AGN emission lines} in the KCWI IFU central spectrum
  (Córdova Rosado et al., in prep.).
\item
  \textbf{X-ray:} no Chandra source at the galaxy position;
  \textbf{0.5--7 keV flux upper limit 1.8×10⁻¹⁵ erg s⁻¹ cm⁻²} (Chandra
  DR2; field obs PI Hughes, ID 16229, 30 ks) → \textbf{L\_X \textless{}
  10⁴³ erg s⁻¹ ≈ 10⁻⁴ × L\_Edd} (the draft uses a 10⁹ M☉ BH; for our
  canonical \textbf{5.2×10⁸} M☉, L\_Edd halves → ≈2×10⁻⁴ L\_Edd, still
  order 10⁻⁴ --- the quiescence conclusion is unchanged).
\item
  \textbf{Radio-quiet:} below the \textbf{LOFAR DR3} detection limit.
\item
  → The central gravitational perturbation is from a \textbf{quiescent
  (non-accreting) SMBH} --- the headline novelty (a dynamical-equivalent
  mass for a quiescent BH at cosmological distance). \textbf{G2 is now
  partly internal}, but the X-ray flux limit and the Chandra/LOFAR
  program details should still be cross-checked against the actual
  Chandra DR2 / LOFAR DR3 sources before we assert them in \emph{our}
  paper.
\end{itemize}

\begin{center}\rule{0.5\linewidth}{0.5pt}\end{center}

\hypertarget{key-citations-verified-in-notebook-10}{%
\section{Key citations (verified in notebook
10)}\label{key-citations-verified-in-notebook-10}}

\begin{itemize}
\tightlist
\item
  \textbf{ppxf:} Cappellari \& Emsellem (2004); Cappellari (2017, MNRAS
  466, 798); Cappellari (2023).
\item
  \textbf{σ\_e definition:} \textbf{Gültekin et al.~(2009)} eq. 1 (the
  luminosity-weighted second-moment integral --- verified 2026-06-12);
  co-add-then-fit method \textbf{Cappellari et al.~(2006, MNRAS 366,
  1126, SAURON IV §2.3)} {[}NB: that paper's eq. 1 is the \emph{aperture
  correction}, NOT σ\_e --- do not cite ``Cappellari 2006 eq. 1'' for
  σ\_e{]}; the discrete spaxel/bin sum (§7) is \textbf{Cappellari et
  al.~(2013, ATLAS3D XV, MNRAS 432, 1709) eq. 29} (= the discretised
  Gültekin 2009 eq. 1; eq. 29 also has an inclination/projection term
  that our observed-plane sum omits) {[}verified 2026-06-12{]};
  \textbf{Kormendy \& Ho (2013, ARA\&A 51, 511)} eq. 3 (scatter 0.29
  dex); \textbf{Greene, Strader \& Ho (2020, ARA\&A 58, 257)} fig.~5
  \textbf{{[}⚠ AUDIT FLAG --- see below{]}}; Saglia et al.~(2016).
\item
  \textbf{z=0 IFS calibration:} Cappellari et al.~(2013, ATLAS3D XV \&
  XX); Krajnović et al.~(2013, XVII); Zhu et al.~(2024, MaNGA DynPop
  III).
\item
  \textbf{Mass--size / analogs:} van der Wel et al.~(2014); Mowla et
  al.~(2019); Ge et al.~(2021); Sonnenfeld et al.~(2015, SL2S V)
  \textbf{{[}⚠ AUDIT FLAG --- see below{]}}.
\item
  \textbf{Vacuum↔air:} Ciddor (1996). \textbf{Wide-window precedent:}
  Bezanson et al.~(2018); van der Wel et al.~(2021).
\item
  \textbf{Source-z catalog:} AGEL DR2 (Carleton et al., in prep).
  \textbf{Cluster:} Hilton et al.~(2021, ACT DR5).
\end{itemize}

\hypertarget{citation-audit-referee-challenges-2026-06-08}{%
\section{Citation audit + referee challenges
(2026-06-08)}\label{citation-audit-referee-challenges-2026-06-08}}

\textbf{Citation audit} (vs the user's Zotero library; ``not in
library'' ≠ fabricated --- most are real foundational papers simply not
yet added to Zotero, flagged so the .bib can be completed):

\begin{itemize}
\tightlist
\item
  \textbf{Verified in-library, metadata exact:} Cappellari (2017, MNRAS
  466, 798) ✓; Kormendy \& Ho (2013, ARA\&A 51, 511) ✓; Greene, Strader
  \& Ho (2020, ARA\&A 58, 257) ✓ \emph{(metadata only --- see flag)};
  Hilton et al.~(2021, ApJS 253, 3 = ACT SZ-cluster catalog) ✓;
  Cappellari (2023, MNRAS 526, 3273) ✓. Also in-library and relevant:
  \textbf{McConnell \& Ma (2013)} --- a massive-BH M•--σ / M•--M\_bulge
  anchor.
\item
  \textbf{⚠ FLAG 1 --- Greene, Strader \& Ho (2020, ARA\&A 58, 257):}
  bibliographic metadata is CORRECT, but this review is titled
  \textbf{``Intermediate-Mass Black Holes.''} It does provide M•--σ /
  M•--M★ relations (legitimately citable as a \emph{comparison}
  relation, as Fig. 4 uses it), but for a \textasciitilde10⁹ M☉ SMBH the
  \emph{primary} anchors should be KH13 / McConnell \& Ma 2013 /
  Saglia+2016, with Greene+20 as one of the overplotted comparison
  relations. \textbf{ACTION: verify the ``fig.~5'' number} (the
  scaling-relation figure in Greene+20 may not be fig.~5) and ensure the
  text frames it as a comparison, not the primary massive-BH relation.
  \emph{(Not a fabrication --- a framing/figure-number check.)}
\item
  \textbf{⚠ FLAG 2 --- Sonnenfeld ``SL2S V'' (2015):} the library
  contains SL2S \textbf{IV} (Sonnenfeld et al. \textbf{2013}, ApJ 777,
  98), not paper \textbf{V} (2015, ApJ 800, 94). \textbf{ACTION: decide
  which you mean} and either correct to ``IV (2013)'' or add the 2015 V
  paper to Zotero. Do not let the 2013 record stand in for a 2015
  citation.
\item
  \textbf{⚠ FLAG 3 --- Bezanson et al.~(2018):} no Bezanson-first-author
  2018 item in the library. If you mean the LEGA-C velocity-dispersion
  paper (Bezanson et al.~2018, ApJ 858, 60), \textbf{add it} / verify.
\item
  \textbf{To add to Zotero before the .bib is complete} (real papers,
  just absent; metadata NOT independently verified here --- confirm on
  add): Cappellari \& Emsellem (2004), Cappellari et al. (2006, SAURON
  IV), Gültekin et al.~(2009), Saglia et al.~(2016), Cappellari et
  al.~(2013, ATLAS3D), Krajnović et al.~(2013), Zhu et al.~(2024, MaNGA
  DynPop III), van der Wel et al.~(2014), Mowla et al.~(2019), Ge et
  al.~(2021), Ciddor (1996), van der Wel et al.~(2021, LEGA-C DR3).
\end{itemize}

\textbf{Referee challenges to the current doc} (anticipate / address in
drafting):

\begin{enumerate}
\def\labelenumi{\arabic{enumi}.}
\tightlist
\item
  \textbf{R\_e method family (was ``non-converged''; largely ADDRESSED
  2026-06-11).} Headline R\_e = 2.097″ (best-mask raw CoG, r\_max=6″) is
  now \textbf{photutils-validated to ±0.002″}
  (\BreakablePath{validate\_\discretionary{}{}{}Re\_\discretionary{}{}{}photutils.\discretionary{}{}{}py}:
  independent \BreakablePath{RadialProfile} integration), and the raw
  CoG sits at the \textbf{top} of a tight best-mask method family ---
  sky-sub CoG 1.922″, Sérsic r\_eff 1.897″ → method systematic
  \textbf{±0.100″} carried explicitly. The raw-CoG-is-top-of-family
  caveat remains (no sky pedestal subtraction), but it is now bounded
  and documented, not a flagged unknown. Knock-on: R\_e sets the σ\_e
  aperture and (weakly) M★; both effects are inside their systematic
  budgets.
\item
  \textbf{R\_e-source budget term (RESOLVED 2026-06-11).} Now uses the
  \textbf{best-mask CoG light family only} \{1.912/2.097/2.281″\} →
  ±5.13 --- the referee-preferred light-family treatment. CaHK+G 2.90″
  (a different I-map definition) and the full grid (±9.98) are reported
  as cross-checks, NOT folded, so the budget no longer conflates
  aperture-uncertainty with σ(R)-physics.
\item
  \textbf{Rising outer σ(r) attributed to arc dilution.} The claim
  ``σ(r) should decrease; any rise is the z=1.302 arc'' is the right
  physical call, but it is \emph{also} what lets us drop the high-σ
  outer bins; make the arc-contamination evidence explicit (EW /
  continuum-dilution diagnostic, nb07e) so it does not read as motivated
  reasoning.
\item
  \textbf{He I 3819 source-emission mask} rests on a 3-pixel residual
  cluster --- defensible (consistent across reductions, matches
  z=1.302), but small; a referee may probe it. The M8/M9/M10 audit trail
  covers this.
\item
  \textbf{σ\_inst = 12.6 km/s vs σ\_e ≈ 270:} well-resolved, not a
  challenge --- but state the LSF treatment.
\end{enumerate}

\hypertarget{flagged-open-items-carry-into-drafting}{%
\section{Flagged open items (carry into
drafting)}\label{flagged-open-items-carry-into-drafting}}

\begin{enumerate}
\def\labelenumi{\arabic{enumi}.}
\tightlist
\item
  \sout{R\_e-source systematic not folded in} \textbf{RESOLVED
  2026-06-08 (M12, nb15):} re-measured at the wide window (±6.13) and
  folded into the budget → headline ±11.77 → \textbf{±13.27}.
\item
  Reduction-pass ±3.45 rests on only 2 reductions --- refine if a 3rd
  lands.
\item
  \sout{Hδ may still need targeted masking} \textbf{RESOLVED 2026-06-08
  (M12, nb16):} decision = \textbf{keep unmasked} (Hδ is well-fit, not a
  contaminant; masking it would discard LOSVD information). TODO in
  \BreakablePath{bootstrap\_\discretionary{}{}{}ppxf.\discretionary{}{}{}py}
  closed. 3b. \textbf{NEW (A3c, nb14):} raw CoG R\_e is
  r\_max-non-convergent; headline 2.305″ is the top of a 2.1--2.5″
  family (\textbf{R\_e method systematic ±0.08″}), not folded into
  σ\_e/M★ --- flagged for a decision.
\item
  Asymmetric-error tables: older 4-corner/M10 tables print −17.92/+17.65
  (pre-M11), M11 was −11.98/+11.57 (±11.77); \textbf{current (M12) is
  −13.45/+13.10 (±13.27). Quote the M12 value} (canonical:
  \BreakablePath{results/\discretionary{}{}{}PAPER\_\discretionary{}{}{}VALUES.\discretionary{}{}{}json}).
\item
  No X-ray / quiescent data in repo (G2). 6. No lens-model content in
  repo (G1).
\item
  ``We do not account for peculiar velocities'' needs to be written in
  (G3 --- true, just add it).
\item
  Source-z (1.302) error bar + its cache are a \BreakablePath{[TBD]}
  placeholder.
\item
  K409 (Aug-30) PI and per-night DIMM seeing still TBD for
  acknowledgements.
\end{enumerate}

\begin{center}\rule{0.5\linewidth}{0.5pt}\end{center}

\hypertarget{internal-appendix-ux3c3_e-measurement-audit-trail-m8m12-not-for-the-manuscript}{%
\section{Internal appendix --- σ\_e measurement audit trail (M8→M12) ---
NOT for the
manuscript}\label{internal-appendix-ux3c3_e-measurement-audit-trail-m8m12-not-for-the-manuscript}}

\begin{quote}
\textbf{Superseded by ``best mask throughout'' (2026-06-11).} The
numbers in THIS appendix (σ\_e=269.62, total ±13.27, R\_e-source ±6.13,
R\_e=2.305″) are the \textbf{M12 (2026-06-08) state}. They were the
headline until the best-mask cascade adopted R\_e=2.097″ (best-mask CoG,
photutils-validated), which moved σ\_e → \textbf{267.31 ± 12.79} and
re-derived R\_e-source → \textbf{±5.13} (best-mask CoG family;
\BreakablePath{run\_\discretionary{}{}{}sigma\_\discretionary{}{}{}e\_\discretionary{}{}{}Re\_\discretionary{}{}{}grid.\discretionary{}{}{}py}).
The current state is §2.1.2 / §2.4.3 / §3.1 above. This trail is kept
verbatim as M12 provenance; do not quote its numbers as current.
\end{quote}

\emph{This is the iterative history of how the final §2.4.3 masking +
systematic-component set was reached. The Methods section should present
the \textbf{final} procedure (§2.4) and the budget (§3.1), not this
trail. Kept here as internal provenance / referee-rebuttal backup. Each
step is also a \BreakablePath{TESTS} row with its own cache.}

\begin{itemize}
\tightlist
\item
  \textbf{M8 --- He I 3819 (added as a source-emission mask).} A 3-pixel
  positive-residual cluster at def-rest 5244--5248 Å = source-rest
  3818--3820 Å matches He I 3819.6 from the z=1.302 source; consistent
  across both reductions → astrophysical, not a CR. Masked. σ\_e (NEW)
  268.98 → 271.87.
\item
  \textbf{M9 --- 5193--5204 Å bump (deliberately NOT masked).} A +5--7\%
  feature; checked against the arc spectrum (no source counterpart) and
  the noise spectrum (no sky counterpart) → it is the \textbf{Mg b LOSVD
  wing} (signal). Smoke-masking it drops σ\_e by 7 km/s → \textbf{DO NOT
  mask.} (Same ``it's signal, not contamination'' logic later applied to
  Hδ in M12.)
\item
  \textbf{M10 --- full sky-line audit.} Flag every band with cube
  \BreakablePath{noise\_\discretionary{}{}{}sky} \textgreater{} 2.5×
  median across the fit window, and cross-check each against the arc
  spectrum to confirm no source counterpart → 9 OH/sky bands added to
  \BreakablePath{BAD\_\discretionary{}{}{}PIXELS\_\discretionary{}{}{}REST}
  (26 → 35 entries). σ\_e (NEW) 271.87 → 269.62.
\item
  \textbf{M11 --- cube-matched systematic re-derivation.} Re-ran the
  I-shape / mask-weight / fit-window sweeps on the NEW cube + M10 masks
  (replacing carried OLD-cube values). The fit-window term collapsed ±15
  → ±3.82 on the cleaned cube; total sys ±17.16 → ±10.81 (total ±11.77).
\item
  \textbf{M12 (2026-06-08) --- R\_e-source fold-in + two cross-checks.}
  Folded the wide-window R\_e-source systematic (±6.13) into the budget
  → sys ±12.43, \textbf{total ±13.27}. Added the
  arc-mask-\textbf{definition} cross-check (±4.58; overlaps the
  F200-mask term, not double-counted) and resolved the Hδ open item
  (\textbf{keep unmasked} --- well-fit, masking it discards LOSVD
  information). Central σ\_e 269.62 unchanged across M10→M12; only the
  error grew.
\end{itemize}

\textbf{Net effect of the trail:} central σ\_e settled at 269.62 km/s by
M10 and has not moved since; the error bar evolved ±17.78 (pre-M11) →
±11.77 (M11) → ±13.27 (M12). The mask set (He I + 35-entry sky/CR +
no-Balmer) and the 7-component budget that the final §2.4 procedure uses
are the M12 result.

\end{document}
